\chapter{Depth, regularity, duality}

In this chapter, all rings are assumed commutative, algebras are associative, commutative and unital, and algebra homomorphisms are unital. If A is a local ring, one denotes \(m_{A}\) its maximal ideal, and \(k_{A}\) its residue field. If p is a prime ideal of a ring A, one denotes \(\kappa(\mathfrak{p})\) the residue field of the local ring \(A_{p}\) ; one identifies it with the field of fractions of the integral domain A/p. One denotes \(\overline{Z}\) the subset \(Z \cup \{-\infty, +\infty\}\) of \(\overline{R}\) . One therefore has, for all \(n \in Z\) , the relations \(-\infty < n < +\infty\) and \(n + \infty = \infty + n = \infty + \infty = \infty\) , \(n - \infty = -\infty + n = -\infty - \infty = -\infty\) .

\section{Depth}

\subsection{Homological definition of depth}

Let A be a ring, J an ideal of A, and M an A-module.

DEFINITION 1. The depth of M with respect to J is called and denoted \(\mathrm{prof}_{\mathrm{A}}(\mathrm{J};\mathrm{M})\) , or \(\mathrm{prof}(\mathrm{J};\mathrm{M})\) , the lower bound in \(\mathbf{N}\cup\{+\infty\}\) of the set of integers \(n\) such that \(\mathrm{Ext}_{\mathrm{A}}^{n}(\mathrm{A}/\mathrm{J},\mathrm{M})\) is nonzero.

When the ring A is local, one simply calls the depth of M and denotes \(\text{prof}_{\mathrm{A}}(\mathrm{M})\) or \(\text{prof}(\mathrm{M})\) the depth of M with respect to the maximal ideal \(m_{A}\) of A; one calls the depth of the local ring A the depth of the A-module A.

\(\operatorname{Si} \operatorname{prof}_{\mathrm{A}}(\mathrm{J}; \mathrm{M}) = +\infty\) , one has \(\operatorname{Ext}_{\mathrm{A}}^{i}(\mathrm{A}/\mathrm{J}, \mathrm{M}) = 0\) for all \(i\) . \(\operatorname{Si} \operatorname{prof}_{\mathrm{A}}(\mathrm{J}; \mathrm{M}) = r < +\infty\) , one has \(\operatorname{Ext}_{\mathrm{A}}^{i}(\mathrm{A}/\mathrm{J}, \mathrm{M}) = 0\) for \(i < r\) and \(\operatorname{Ext}_{\mathrm{A}}^{r}(\mathrm{A}/\mathrm{J}, \mathrm{M}) \neq 0\) .

Remarks.--- 1) Suppose that the A-module M is finitely generated and that one has M = JM, that is, Supp(M) ∩ V(J) = ∅ (II, § 4, no. 4, cor. of prop. 18). In this case, prof\_A(J;M) is equal to +∞: indeed, the ideal Ann(M) + J is then equal to A (loc. cit.) and is contained in the annihilator of Ext\_A\^{}i(A/J;M) for all i. One will see below (no. 3, cor. 1 of th. 1) that conversely if the ideal J is finitely generated, prof\_A(J;M) = +∞ implies M = JM.

\begin{enumerate}
\def\labelenumi{\arabic{enumi})}
\setcounter{enumi}{1}
\tightlist
\item
  For \(\mathrm{prof}_{\mathbf{A}}(\mathbf{J};\mathbf{M})\) to be zero, it is necessary and sufficient that \(\mathrm{Hom}_{\mathbf{A}}(\mathbf{A} / \mathbf{J},\mathbf{M})\) be nonzero, that is, that M possess a nonzero element annihilated by J; this is in particular the case when one has \(\mathrm{Ass(M)}\cap \mathrm{V(J)}\neq \emptyset\) . If the ring A is noetherian and the A-module M is finitely generated, the following conditions are equivalent (IV, § 1, no. 4, prop. 8):
\end{enumerate}

\begin{enumerate}
\def\labelenumi{\alph{enumi})}
\item
  \(\mathrm{prof}_{\mathrm{A}}(\mathrm{J};\mathrm{M}) = 0\)
\item
  for every \(x \in J\) , the homothety \(x_{\mathbb{M}}\) is not injective;
\end{enumerate}

\[
\mathrm{c)} \text {   we have   } \operatorname{Ass} (\mathrm{M}) \cap \mathrm{V} (\mathrm{J}) \neq \varnothing .
\]

\begin{enumerate}
\def\labelenumi{\arabic{enumi})}
\setcounter{enumi}{2}
\item
  According to remark 2, for a local ring \(\Lambda\) to be of depth zero, it is necessary and sufficient that there exist a nonzero element \(x\) of A such that \(\mathfrak{m}_{\mathrm{A}}x = 0\) . If A is not a field, an element \(x\) such that \(\mathfrak{m}_{\mathrm{A}}x = 0\) is not invertible, hence belongs to \(\mathfrak{m}_{\mathrm{A}}\) and consequently satisfies \(x^{2} = 0\) . Thus a reduced local ring of dimension \(\geqslant 1\) is of depth \(\geqslant 1\) .
\item
  Let \((\mathbf{M}_{\iota})_{\iota \in I}\) be a family of A-modules. According to A, X, p. 89, prop. 7, one has \(\operatorname{prof}(\mathbf{J};\prod_{\iota \in I}\mathbf{M}_{\iota}) = \inf_{\iota \in I}\operatorname{prof}(\mathbf{J};\mathbf{M}_{\iota})\) .
\end{enumerate}

PROPOSITION 1.--- Let A be a ring, J an ideal of A, and 0 → M' → M → M'' → 0 an exact sequence of A-modules. Set

\[
p ^ {\prime} = \operatorname{prof} (\mathrm{J}; \mathrm{M} ^ {\prime}), p = \operatorname{prof} (\mathrm{J}; \mathrm{M}), p ^ {\prime \prime} = \operatorname{prof} (\mathrm{J}; \mathrm{M} ^ {\prime \prime}).
\]

One is then in one of the following three cases, which are mutually exclusive:

\[
p ^ {\prime} = p \leqslant p ^ {\prime \prime}, \quad p = p ^ {\prime \prime} <   p ^ {\prime}, \quad p ^ {\prime \prime} = p ^ {\prime} - 1 <   p.
\]

Consider the exact sequence of extension modules associated with A/J and with the above exact sequence (A, X, p. 92, th. 2). Excluding the case \(p = p' = p'' = +\infty\) ; there then exists in this sequence a first nonzero module, and the following module is also nonzero. This gives the following three possibilities:

\begin{enumerate}
\def\labelenumi{\alph{enumi})}
\item
  The first nonzero module is \(\operatorname{Ext}_{\Lambda}^{p'}(A/J, M')\) . One then has \(p' = p \leqslant p''\) .
\item
  The first nonzero module is \(\operatorname{Ext}_{\Lambda}^{p}(\mathrm{A} / \mathrm{J},\mathrm{M})\) . One then has \(p = p'' < p'\) .
\item
  The first nonzero module is \(\mathrm{Ext}_{\mathrm{A}}^{p''}(\mathrm{A} / \mathrm{J},\mathrm{M}'')\) . One then has \(p'' + 1 = p' \leqslant p\) .
\end{enumerate}

Remark 5.--- Suppose that one has \(p = p'\) and that the injection \(u: M' \to M\) which occurs in the exact sequence of prop. 1 belongs to J Hom \(_A(M', M)\) . One \(a\) then \(p'' = p - 1\) . Indeed, the hypothesis implies that the map Ext \(_A^i(1_{A/J}, u)\) is zero for every integer \(i\) ; this excludes the case \(a\) ) considered above.

PROPOSITION 2.--- Let A be a ring, J an ideal of A, M an A-module, and N an A-module annihilated by a power of J. Then \(\mathrm{Ext}_{\mathrm{A}}^{i}(\mathrm{N},\mathrm{M}) = 0\) for every integer \(i < \mathrm{prof}_{\mathrm{A}}(\mathrm{J};\mathrm{M})\) .

Suppose first that JN = 0 and argue by induction on the integer \(i < \mathrm{prof}_{\mathrm{A}}(\mathrm{J};\mathrm{M})\) . The assertion is evident for \(i < 0\) . Consider N as an (A/J)-module and choose an exact sequence of (A/J)-modules

\[
0 \rightarrow \mathrm{K} \longrightarrow (\mathrm{A} / \mathrm{J}) ^ {(\mathrm{I})} \longrightarrow \mathrm{N} \rightarrow 0.
\]

From this we deduce an exact sequence of extension modules

\[
\operatorname{Ext} _ {\mathrm{A}} ^ {i - 1} (\mathrm{K}, \mathrm{M}) \longrightarrow \operatorname{Ext} _ {\mathrm{A}} ^ {i} (\mathrm{N}, \mathrm{M}) \longrightarrow \operatorname{Ext} _ {\mathrm{A}} ^ {i} ((\mathrm{A} / \mathrm{J}) ^ {(1)}, \mathrm{M}).
\]

The A-module \(\mathrm{Ext}_{\mathrm{A}}^{i-1}(\mathrm{K},\mathrm{M})\) is zero by the induction hypothesis, and the A-module \(\mathrm{Ext}_{\mathrm{A}}^{i}((\mathrm{A}/\mathrm{J})^{(1)},\mathrm{M})\) is isomorphic to \(\mathrm{Ext}_{\mathrm{A}}^{i}(\mathrm{A}/\mathrm{J},\mathrm{M})^{\mathrm{I}}\) (A, X, p.~89, prop. 7), which is zero by definition of the depth. Consequently we have \(\mathrm{Ext}_{\mathrm{A}}^{i}(\mathrm{N},\mathrm{M}) = 0\) .

Let us pass to the general case and argue by induction on the smallest integer m \textgreater{} 0 such that \(J^{m}N = 0\) . We have just treated the case m = 1. Suppose m \textgreater{} 1 and let \(i < \text{prof}_{\text{A}}(\text{J}; \text{M})\) be an integer. Consider the exact sequence

\[
\operatorname{Ext} _ {\Lambda} ^ {i} (\mathrm{N/JN,M}) \longrightarrow \operatorname{Ext} _ {\Lambda} ^ {i} (\mathrm{N,M}) \longrightarrow \operatorname{Ext} _ {\Lambda} ^ {i} (\mathrm{JN,M})
\]

deduced from the exact sequence \(0 \rightarrow JN \rightarrow N \rightarrow N/JN \rightarrow 0\) . The two extreme modules are zero by the induction hypothesis, since N/JN and JN are annihilated by \(J^{m-1}\) . One therefore has \(\operatorname{Ext}_{\Lambda}^{i}(N,M) = 0\) , which is what we wanted to prove.

COROLLARY 1.--- Let m be an integer \textgreater0 and let J' be an ideal of A containing J\^{}m. Then \(\operatorname{prof}_{\mathrm{A}}(\mathrm{J};\mathrm{M}) \leqslant \operatorname{prof}_{\mathrm{A}}(\mathrm{J}';\mathrm{M})\) .

Indeed \(J^{m}\) annihilates the A-module \(A/J'\) , hence \(\operatorname{Ext}_{A}^{i}(A/J', M)\) is zero for every integer \(i < \operatorname{prof}_{A}(J; M)\) (prop. 2).

COROLLARY 2.--- Suppose the ideal J is finitely generated, and let J' be an ideal of A such that V(J) ⊃ V(J').

\begin{enumerate}
\def\labelenumi{\alph{enumi})}
\item
  One \(a \operatorname{prof}_{\mathrm{A}}(\mathrm{J}; \mathrm{M}) \leqslant \operatorname{prof}_{\mathrm{A}}(\mathrm{J}'; \mathrm{M})\) .
\item
  If the ideal J' is finitely generated and if V(J) = V(J'), then we have \(\text{prof}_{A}(J; M) = \text{prof}_{A}(J'; M)\) .
\end{enumerate}

By II, § 4, no. 3, cor. 2 of prop. 11 and § 2, no. 6, prop. 15, there exists an integer m \textgreater{} 0 such that \(J^{m} \subset J'\) . Assertion a) therefore follows from cor. 1, and assertion b) is deduced from it.

Cor. 2 may fail when the ideal J is not finitely generated (exercise 2).

\subsection{Depth and acyclicity}

PROPOSITION 3.--- Let A be a ring, C a bounded-below complex of A-modules, and p an integer. Suppose that for every pair of integers \((m,n)\) with \(m \geqslant n \geqslant p\) , the depth of the A-module \(C_{m}\) relative to the annihilator of \(\mathrm{H}_{n}(\mathrm{C})\) is \textgreater m - n. We then have \(\mathrm{H}_{n}(\mathrm{C}) = 0\) for \(n \geqslant p\) .

Since C is bounded on the left, \(\mathrm{H}_{n}(\mathrm{C})\) is zero for n sufficiently large. If the conclusion were false, there would exist an integer \(q \geqslant p\) such that \(\mathrm{H}_{n}(\mathrm{C}) = 0\) for n \textgreater{} q and \(\mathrm{H}_{q}(\mathrm{C}) \neq 0\) . Let J denote the annihilator of \(\mathrm{H}_{q}(\mathrm{C})\) ; we then have \(\operatorname{prof}_{\mathrm{A}}(\mathrm{J}; \mathrm{H}_{q}(\mathrm{C})) = 0\) . Moreover, since \(\mathrm{Z}_{q}(\mathrm{C})\) is a submodule of \(C_{q}\) , and since by hypothesis \(\operatorname{prof}_{\mathrm{A}}(\mathrm{J}; \mathrm{C}_{q}) > q - q = 0\) , one has \(\operatorname{prof}_{\mathrm{A}}(\mathrm{J}; \mathrm{Z}_{q}(\mathrm{C})) > 0\) . We then deduce from the exact sequence

\[
0 \to \mathrm{B} _ {q} (\mathrm{C}) \to \mathrm{Z} _ {q} (\mathrm{C}) \to \mathrm{H} _ {q} (\mathrm{C}) \to 0
\]

the equality \(\mathrm{prof}_{\mathrm{A}}(\mathrm{J};\mathrm{B}_q(\mathrm{C})) = 1\) (n° 1, prop. 1). By the definition of \(q\) , \(\mathrm{B}_n(\mathrm{C})\) is equal to \(\mathrm{Z}_n(\mathrm{C})\) for every integer \(n > q\) . From the canonical exact sequences

\[
0 \to \mathrm{B} _ {n} (\mathrm{C}) \to \mathrm{C} _ {n} \to \mathrm{B} _ {n - 1} (\mathrm{C}) \to 0 \quad (n > q)
\]

and from the hypothesis \(\operatorname{prof}_{\mathrm{A}}(\mathrm{J};\mathrm{C}_{n}) > n - q\) , we obtain by induction the equality \(\operatorname{prof}_{\mathrm{A}}(\mathrm{J};\mathrm{B}_{n}(\mathrm{C})) = n - q + 1\) for every \(n \geqslant q\) (loc. cit.). But this is absurd, since \(\mathrm{B}_{n}(\mathrm{C})\) is zero for \(n\) sufficiently large.

COROLLARY 1.--- Let A be a ring, J an ideal of A, C a bounded-below complex of A-modules, and p an integer. Suppose that \(\mathrm{JH}_{m}(\mathrm{C}) = 0\) and \(\operatorname{prof}_{\mathrm{A}}(\mathrm{J}; \mathrm{C}_{m}) > m - p\) for \(m \geqslant p\) . One then has \(\mathrm{H}_{n}(\mathrm{C}) = 0\) for \(n \geqslant p\) .

Indeed, for \(n \geqslant p\) the annihilator \(J_n\) of \(H_n(C)\) contains \(J\) ; therefore we have \(\mathrm{prof}_{\mathrm{A}}(\mathrm{J}_{n};\mathrm{C}_{m}) \geqslant \mathrm{prof}_{\mathrm{A}}(\mathrm{J};\mathrm{C}_{m})\) ( \(n^{\circ}\) 1, cor. 1 of prop. 2), so that the hypothesis of the proposition is satisfied.

COROLLARY 2.-- Let A be a local ring, C a bounded-below complex of A-modules, p an integer. Suppose that for \(m \geqslant p\) , \(\mathrm{H}_{m}(\mathrm{C})\) is of finite length and \(\mathrm{C}_{m}\) of depth \(>m-p\) . One then has \(\mathrm{H}_{n}(\mathrm{C}) = 0\) for \(n \geqslant p\) .

The A-module \(\bigoplus_{m\geqslant p}\mathrm{H}_{m}(\mathrm{C})\) is of finite length. Let J denote its annihilator; by A, VIII, § 1, no. 3, corollary, the ring A/J is Artinian, hence J contains a power of the maximal ideal of A (A, VIII, § 10, no. 1, th. 1). Consequently one has \(\mathrm{prof}_A(\mathrm{J};\mathrm{C}_m)\geqslant \mathrm{prof}(\mathrm{C}_m) > m - p\) for \(m\geqslant p\) (no. 1, cor. 1 of prop. 2), so that cor. 1 can be applied.

\subsection{Depth and Koszul complex}

Let A be a ring, M an A-module, \(\mathbf{x} = (x_i)_{i \in I}\) a family of elements of A. Denote by \(u : \mathrm{A}^{(1)} \to \mathrm{A}\) the linear form such that \(u(c_i) = x_i\) for every \(i \in I\) , and by \(\mathbf{K}^{\bullet}(\mathbf{x}, \mathrm{M})\) the complex \(\mathbf{K}_{\mathrm{A}}^{\bullet}(u, \mathrm{M})\) associated with u (A, X, p.~147). One has \(\mathbf{K}^p(\mathbf{x}, \mathrm{M}) = 0\) for p \textless{} 0 ; for \(p \geqslant 0\) the A-module \(\mathbf{K}^p(\mathbf{x}, \mathrm{M}) = \operatorname{Hom}_{\Lambda} (\boldsymbol{\Lambda}^p(\mathrm{A}^{(1)}), \mathrm{M})\) is canonically identified with the A-module \(C_I^p(\mathrm{M})\) formed by the alternating maps from \(I^p\) to M (A, X, p.~153), the differential \(\partial^p : \mathbf{K}^p(\mathbf{x}, \mathrm{M}) \longrightarrow \mathbf{K}^{p+1}(\mathbf{x}, \mathrm{M})\) being given by the formula

\[
\left(\partial^ {p} \boldsymbol {m}\right) \left(\alpha_ {1}, \dots , \alpha_ {p + 1}\right) = \sum_ {j = 1} ^ {p + 1} (- 1) ^ {j + 1} x _ {\alpha_ {j}} \boldsymbol {m} \left(\alpha_ {1}, \dots , \alpha_ {j - 1}, \alpha_ {j + 1}, \dots , \alpha_ {p + 1}\right)
\]

for \(\boldsymbol{m} \in \mathbf{K}^{p}(\mathbf{x}, \mathrm{M})\) and \((\alpha_{1}, \ldots, \alpha_{p+1}) \in \mathrm{I}^{p+1}\) (A, X, p. 154, formula (12)). It follows in particular that the complex \(\mathbf{K}^{\bullet}(\mathbf{x}, \mathrm{M})\) depends only on the structure of M as a Z-module and on the endomorphisms \((x_{i})_{\mathrm{M}}\) .

One denotes by \(\mathrm{H}^{\bullet}(\mathbf{x},\mathrm{M})\) the homology of the complex \(\mathbf{K}^{\bullet}(\mathbf{x},\mathrm{M})\) . The A-module \(\mathrm{H}^{0}(\mathbf{x},\mathrm{M})\) is identified with \(\mathrm{Hom}_{\Lambda}(\mathrm{A}/\mathrm{J},\mathrm{M})\) , where J is the ideal of A generated by the \(x_{i}\) (A, X, p.~147, lemma 1).

Let \((\mathrm{M}_{\alpha})_{\alpha\in\mathrm{K}}\) be a family of A-modules, and M its product; the complex \(\mathbf{K}^{\bullet}(\mathbf{x},\mathrm{M})\) is canonically isomorphic to the product complex of the \(\mathbf{K}^{\bullet}(\mathbf{x},\mathrm{M}_{\alpha})\) , so that for each integer s the A-module \(\mathrm{H}^{s}(\mathbf{x},\mathrm{M})\) is identified with the product of the \(\mathrm{H}^{s}(\mathbf{x},\mathrm{M}_{\alpha})\) (A, X, p. 28, prop. 1).

THEOREM 1.-- Let A be a ring, J an ideal of A, \(\mathbf{x} = (x_i)_{i \in I}\) a generating family of J, M an A-module. The depth of M relative to J is the lower bound (in \(\mathbf{N} \cup \{+\infty\}\) ) of the integers \(n\) such that \(\Pi^n(\mathbf{x}, \mathbf{M}) \neq 0\) .

Set \(p = \text{prof}_{\Lambda}(\text{J}; \text{M})\) . Consider the complex \(\mathbf{K}^{\bullet}(\mathbf{x}, \text{M})\) . Its homology is annihilated by J ( \(\Lambda\) , X, p.~148, cor. 2), and the depth relative to J of each of the modules \(\mathbf{K}^{i}(\mathbf{x}, \text{M})\) is equal to \(p\) or to \(+\infty\) (no. 1, remark 4). It then follows from cor. 1 of no. 2 that one has \(\text{H}^{i}(\text{x}, \text{M}) = 0\) for \(i < p\) . It remains to prove that \(\text{H}^{p}(\text{x}, \text{M})\) is nonzero when \(p < +\infty\) .

The case \(p = 0\) being evident, suppose \(0 < p < +\infty\) and \(\mathrm{H}^p(\mathbf{x},\mathrm{M}) = 0\) . Let L be a free resolution of the A-module \(\mathrm{A / J}\) ; denote by C the complex \(\mathrm{Homgr}_{\mathrm{A}}(\mathrm{L},\mathrm{M})\) . The \(\Lambda\) -module \(\mathrm{H}^i(\mathrm{C})\) is isomorphic to \(\mathrm{Ext}_{\mathrm{A}}^i(\mathrm{A / J},\mathrm{M})\) (A, X, p.~100, th. 1); it is therefore zero for i \textless{} p.~One then has for i \textless{} p canonical exact sequences

\[
0 \rightarrow \mathrm{B} ^ {i} (\mathrm{C}) \rightarrow \mathrm{C} ^ {i} \rightarrow \mathrm{B} ^ {i + 1} (\mathrm{C}) \rightarrow 0.
\]

The A-module \(C^{i}\) is a product of A-modules isomorphic to M ; hence one has \(\mathrm{H}^{s}(\mathbf{x},\mathrm{C}^{i})=0\) for \(s\leqslant p\) . From the preceding exact sequences and A, X, p.~150 one deduces that the connecting homomorphism \(\partial^{s}:\mathrm{H}^{s}(\mathbf{x},\mathrm{B}^{i+1}(\mathrm{C}))\longrightarrow\mathrm{H}^{s+1}(\mathbf{x},\mathrm{B}^{i}(\mathrm{C}))\) is injective for \(s\leqslant p\) and i\textless p ; since \(\mathrm{B}^{0}(\mathrm{C})=0\) , it follows that \(\mathrm{H}^{p-i}(\mathbf{x},\mathrm{B}^{i+1}(\mathrm{C}))\) is zero for i\textless p.~In particular, we have \(\mathrm{II}^{1}(\mathbf{x},\mathrm{B}^{p}(\mathrm{C}))=0\) , so that the exact sequence

\[
0 \to \mathrm{B} ^ {p} (\mathrm{C}) \to \mathrm{Z} ^ {p} (\mathrm{C}) \to \mathrm{H} ^ {p} (\mathrm{C}) \to 0
\]

provides a surjection \(\mathrm{H}^{0}(\mathbf{x},\mathrm{Z}^{p}(\mathrm{C}))\longrightarrow\mathrm{H}^{0}(\mathbf{x},\mathrm{H}^{p}(\mathrm{C}))\) . Since \(\mathrm{H}^{p}(\mathrm{C})\) is isomorphic to \(\operatorname{Ext}_{\mathrm{A}}^{p}(\mathrm{A}/\mathrm{J},\mathrm{M})\) which is nonzero and annihilated by J, we have \(\mathrm{H}^{0}(\mathbf{x},\mathrm{H}^{p}(\mathrm{C}))\neq0\) , whence \(\mathrm{H}^{0}(\mathbf{x},\mathrm{Z}^{p}(\mathrm{C}))\neq0\) and consequently \(\mathrm{H}^{0}(\mathbf{x},\mathrm{C}^{p})\neq0\) . But this implies \(\mathrm{H}^{0}(\mathbf{x},\mathrm{M})\neq0\) , contrary to the hypothesis. We therefore have \(\mathrm{H}^{p}(\mathbf{x},\mathrm{M})\neq0\) , which completes the proof.

COROLLARY 1.--- Suppose the ideal J is finitely generated and JM ≠ M. Then prof \(_{A}\) (J; M) is finite and ≤ Card(I); for it to be equal to Card(I), it is necessary and sufficient that the family x be completely secant for M (A, X, p. 157, def. 2).

Suppose first that 1 is finite, and denote \(r\) its cardinality. The A-module \(\mathrm{H}^r (\mathbf{x},\mathrm{M})\) is canonically isomorphic to \(\mathrm{H}_0(\mathbf{x},\mathrm{M})\) , itself isomorphic to M/JM (A, X, p.~155); the inequality \(\mathrm{prof}_{\Lambda}(\mathrm{J};\mathrm{M})\leqslant r\) It therefore follows from Th. 1. For equality to hold, it is necessary and sufficient that the A-module \(\mathrm{H}^i (\mathbf{x},\mathrm{M})\) be zero for \(i < r\) , which means that the family \(\mathbf{x}\) is completely secant for M (A, X, p.~157).

Based on what precedes, \(\mathrm{prof}_{\mathrm{A}}(\mathrm{J};\mathrm{M})\) is finite; it remains to prove that if the family \(\mathbf{x}\) is completely secant for M, the set I is finite. But the condition \(\mathrm{H}_{1}(\mathbf{x},\mathrm{M}) = 0\) (A, X, p. 157, def. 2) implies that we have an exact sequence

\[
\boldsymbol {\Lambda} ^ {2} (\mathrm{A} ^ {(\mathrm{I})}) \otimes_ {\mathrm{A}} \mathrm{M} \xrightarrow {\partial_ {2}} \mathrm{M} ^ {(\mathrm{I})} \xrightarrow {\partial_ {1}} \mathrm{JM} \to 0,
\]

where the image of \(\partial_{2}\) is contained in \(\mathrm{JM}^{(1)}\) By taking the tensor product with A/J, we deduce an A-linear isomorphism from \((\mathrm{M} / \mathrm{JM})^{(1)}\) on \(\mathrm{JM} / \mathrm{J}^2\mathrm{M}\) . Now this latter module is of finite type, since J and M are; as M/JM is not zero, it follows that the set I is finite.

COROLLARY 2. Let A be a local ring, J a finitely generated ideal of A distinct from A, M a nonzero finitely generated A-module. Set \(r = [\mathrm{J} / \mathfrak{m}_{\mathrm{A}}\mathrm{J} : \kappa_{\mathrm{A}}]\) . One has \(\operatorname{prof}_{\Lambda}(\mathrm{J};\mathrm{M}) \leqslant r\) ; there is equality if and only if J is generated by a completely secant family for M. In this case, for a generating family of J to be completely secant, it is necessary and sufficient that it have r elements.

By Nakayama's lemma, we have \(\mathrm{JM} \neq \mathrm{M}\) , and \(r\) is the minimal number of generators of \(\mathbf{J}\) ; cor. 2 therefore follows from cor. 1.

PROPOSITION 4.--- Let \(\rho: A \to B\) a homomorphism of rings, \(J\) an ideal of \(A\) and \(N\) a B-module. We have the equality \(\operatorname{prof}_{\mathrm{A}}(\mathrm{J}; \mathrm{N}) = \operatorname{prof}_{\mathrm{B}}(\mathrm{JB}; \mathrm{N})\) .

Let \(\mathbf{x} = (x_{i})_{i \in \mathbb{I}}\) a generating family of J; the family \(\rho(\mathbf{x}) = (\rho(x_{i}))_{i \in \mathbb{I}}\) generates JB. By construction the complex \(\mathbf{K}^{\bullet}(\boldsymbol{\rho}(\mathbf{x}), \mathbf{N})\) is equal to \(\mathbf{K}^{\bullet}(\mathbf{x}, \mathbf{N})\) . The proposition therefore follows from th. 1.

COROLLARY. Let A be a local ring, a an ideal of A distinct from A and M an A-module annihilated by a. We have \(\operatorname{prof}_{\mathrm{A}}(\mathrm{M}) = \operatorname{prof}_{\mathrm{A}/\mathfrak{a}}(\mathrm{M})\) .

Let \(\rho: A \to B\) a homomorphism of rings, \(x = (x_i)_{i \in I}\) a finite family of elements of \(\Lambda\) and M an A-module. For every integer \(p\) let us denote \(u^p: B \otimes_\Lambda C_I^p(M) \longrightarrow C_I^p(B \otimes_\Lambda M)\) the B-linear homomorphism which associates to \(b \otimes m\) the alternating map \((\alpha_1, \ldots, \alpha_p) \mapsto b \otimes m(\alpha_1, \ldots, \alpha_p)\) . The family \((u^p)\) defines an isomorphism of complexes

\[
u: \mathrm{B} \otimes_ {\mathrm{A}} \mathbf {K} ^ {\bullet} (\mathbf {x}, \mathrm{M}) \longrightarrow \mathbf {K} ^ {\bullet} (\mathbf {x}, \mathrm{B} \otimes_ {\mathrm{A}} \mathrm{M}).
\]

Consider the canonical homomorphism

\[
\gamma^ {p} (\mathrm{B}, \mathbf {K} ^ {\bullet} (\mathbf {x}, \mathrm{M})): \mathrm{B} \otimes_ {\mathrm{A}} \mathrm{H} ^ {p} (\mathbf {x}, \mathrm{M}) \longrightarrow \mathrm{H} ^ {p} (\mathrm{B} \otimes_ {\mathrm{A}} \mathbf {K} ^ {\bullet} (\mathbf {x}, \mathrm{M}))
\]

(A, X, p. 62); by composition with \(\mathrm{H}^{p}(u)\) , we deduce from it a homomorphism

\[
v ^ {p}: \mathrm{B} \otimes_ {\Lambda} \mathrm{H} ^ {p} (\mathbf {x}, \mathrm{M}) \longrightarrow \mathrm{H} ^ {p} (\mathbf {x}, \mathrm{B} \otimes_ {\Lambda} \mathrm{M}).
\]

Lemma 1.-- If the A-module B is flat, the homomorphism \(v^p\) is bijective for every integer \(p\) .

This follows from A, X, p. 66, cor. 2.

PROPOSITION 5.--- Let A be a ring, J a finitely generated ideal of A and M an A-module. Let Ω denote the set of maximal ideals of A belonging to Supp(M) and containing J. We have

\[
\operatorname{prof} _ {A} (J; M) = \inf _ {\mathfrak {p} \in \operatorname{Spec} (A)} \operatorname{prof} _ {A _ {\mathfrak {p}}} \left(J _ {\mathfrak {p}}; M _ {\mathfrak {p}}\right) = \inf _ {m \in \Omega} \operatorname{prof} _ {\Lambda_ {m}} \left(J _ {m}; M _ {m}\right).
\]

Let \(\mathbf{x} = (x_{i})_{i \in 1}\) a finite generating family of J. Let p be a prime ideal of A; the ideal \(J_{p}\) is generated by the image \(x_{p}\) of the family \((x_{i})\) in \(A_{p}\) . For every \(p \geqslant 0\) , the \(A_{p}\) -module \(\left(\mathrm{H}^{p}(\mathbf{x},\mathrm{M})\right)_{\mathfrak{p}}\) is isomorphic to \(\mathrm{H}^{p}(\mathbf{x}_{\mathfrak{p}},\mathrm{M}_{\mathfrak{p}})\) (lemma 1); by th. 1, we therefore have \(\operatorname{prof}_{\Lambda}(\mathrm{J};\mathrm{M}) \leqslant \inf_{\mathfrak{p} \in \operatorname{Spec}(\mathrm{A})} \operatorname{prof}_{\mathrm{A}_{\mathfrak{p}}}(\mathrm{J}_{\mathfrak{p}}; \mathrm{M}_{\mathfrak{p}}) \leqslant \inf_{\mathfrak{m} \in \Omega} \operatorname{prof}_{\mathrm{A}_{\mathfrak{m}}}(\mathrm{J}_{\mathfrak{m}}; \mathrm{M}_{\mathfrak{m}})\) .

Let \(p\) be an integer strictly less than \(\operatorname{prof}_{\mathrm{A}_{\mathfrak{m}}}\left(\mathrm{JA}_{\mathfrak{m}};\mathrm{M}_{\mathfrak{m}}\right)\) for every \(\mathfrak{m}\in\Omega\) . One then has \(\mathrm{H}^{p}\left(\mathbf{x}_{\mathfrak{m}},\mathrm{M}_{\mathfrak{m}}\right)=0\) for every maximal ideal \(\mathfrak{m}\) of A: this follows from th. 1 if \(\mathfrak{m}\in\Omega\) , from the fact that \(\mathrm{M}_{\mathfrak{m}}=0\) if \(\mathfrak{m}\notin\operatorname{Supp}(\mathrm{M})\) , and from the fact that the ideal \(\mathrm{JA}_{\mathfrak{m}}\) , which annihilates \(\mathrm{H}^{p}\left(\mathbf{x}_{\mathfrak{m}},\mathrm{M}_{\mathfrak{m}}\right)\) (A, X, p. 148, cor. 2), is equal to \(\mathrm{A}_{\mathfrak{m}}\) if \(\mathfrak{m}\notin\mathrm{V}(\mathrm{J})\) . One therefore has \(\left(\mathrm{H}^{p}(\mathbf{x},\mathrm{M})\right)_{\mathfrak{m}}=0\) for every maximal ideal \(\mathfrak{m}\) of A, which implies \(\mathrm{H}^{p}(\mathbf{x},\mathrm{M})=0\) (II, § 3, n° 3, cor. 2 of th. 1). The proposition then follows from th. 1.

PROPOSITION 6.--- Let A be a ring, J a finitely generated ideal of A and M an A-module. Let B be a ring and \(\rho: A \to B\) a ring homomorphism making B into a flat A-module.

\begin{enumerate}
\def\labelenumi{\alph{enumi})}
\item
  One \(a \operatorname{prof}_{\mathrm{A}}(\mathrm{J}; \mathrm{M}) \leqslant \operatorname{prof}_{\mathrm{B}}(\mathrm{JB}; \mathrm{B} \otimes_{\mathrm{A}} \mathrm{M})\) .
\item
  Suppose moreover that every maximal ideal of Supp(M) containing J belongs to the image of the canonical map \(\operatorname{Spec}(B) \to \operatorname{Spec}(A)\) . One then has \(\operatorname{prof}_{\mathrm{A}}(\mathrm{J}; \mathrm{M}) = \operatorname{prof}_{\mathrm{B}}(\mathrm{JB}; \mathrm{B} \otimes_{\mathrm{A}} \mathrm{M})\) . This is the case in particular if the A-module B is faithfully flat.
\end{enumerate}

Assertion a) follows from th. 1 and lemma 1.

Let \(p\) be an integer strictly less than \(\text{prof}_B(\text{JB};\text{B} \otimes_A \text{M})\) , and let \(\mathfrak{m}\) be a maximal ideal of A belonging to \(\text{Supp(M)} \cap \text{V(J)}\) . Let \(\mathbf{x}\) be a finite generating family of the ideal J. Under the hypothesis of b), there exists a prime ideal \(\mathfrak{n}\) of B lying over \(\mathfrak{m}\) , and one has a canonical isomorphism

\[
\mathrm{B} _ {\mathfrak {n}} \otimes_ {\mathrm{A} _ {\mathfrak {m}}} \left(\mathrm{A} _ {\mathfrak {m}} \otimes_ {\Lambda} \mathrm{H} ^ {p} (\mathbf {x}, \mathrm{M})\right) \longrightarrow \mathrm{B} _ {\mathfrak {n}} \otimes_ {\mathrm{B}} \left(\mathrm{B} \otimes_ {\mathrm{A}} \mathrm{H} ^ {p} (\mathbf {x}, \mathrm{M})\right).
\]

Now \(\mathrm{B}\otimes_{\Lambda}\mathrm{H}^{p}(\mathbf{x},\mathrm{M})\) is isomorphic to \(\mathrm{H}^{p}(\mathsf{p}(\mathbf{x}),\mathrm{B}\otimes_{\Lambda}\mathrm{M})\) (lemma 1), hence is zero; moreover \(B_{n}\) is faithfully flat over \(A_{m}\) (I, § 3, n° 5, prop. 9 and II, § 3, n° 4, prop. 14 and remark). One therefore has \(A_{m}\otimes_{\Lambda}\mathrm{H}^{p}(\mathbf{x},\mathrm{M})=0\) and consequently p \textless{} prof \(_{A_{m}}(J_{m};M_{m})\) (lemma 1 and th. 1). The first assertion of b) then follows from prop. 5; the second follows from I, § 3, n° 5, prop. 9).

COROLLARY.--- Let A be a Noetherian ring, J an ideal of A, M a finitely generated A-module, \(\widehat{\Lambda}\) and \(\widehat{M}\) the separated completions of A and M for the J-adic topology. Then one has \(\operatorname{prof}_{\mathbf{A}}(\mathbf{J};\mathbf{M}) = \operatorname{prof}_{\widehat{\mathbf{A}}}(\widehat{\mathbf{J}}\widehat{\mathbf{A}};\widehat{\mathbf{M}})\) .

Indeed, the A-module \(\widehat{\mathsf{A}}\) is flat and the \(\widehat{\mathsf{A}}\) -module \(\widehat{\mathsf{M}}\) is isomorphic to \(\widehat{\mathsf{A}}\otimes_{\mathsf{A}}\mathsf{M}\) (III, § 3, n° 4, th. 3); moreover, every maximal ideal of A containing J belongs to the image of the map \(\operatorname{Spec}(\widehat{\mathsf{A}})\to\operatorname{Spec}(\mathsf{A})\) (loc. cit., prop. 8).

\subsection{Depth and regular sequences}

Let A be a ring, M an A-module. Recall (A, X, p. 158) that a finite sequence \((x_{1},\ldots,x_{r})\) of elements of A is said to be regular for M, or M-regular, if, for \(i=1,\ldots,r\) , the homothety of ratio \(x_{i}\) in the A-module \(\mathrm{M}/(x_{1}\mathrm{M}+\ldots+x_{i-1}\mathrm{M})\) is injective. Let \((x_{1},\ldots,x_{r})\) be an M-regular sequence; for every flat A-module N, the sequence \((x_{1},\ldots,x_{r})\) is \(M\otimes_{A}N\) regular. If \(\rho:A\to B\) is a ring homomorphism making B into a flat A-module, the sequence \((\rho(x_{1}),\ldots,\rho(x_{r}))\) is regular for the B-module \(B\otimes_{A}M\) . In particular, for every prime ideal p of A, the image in \(A_{p}\) of the sequence \((x_{1},\ldots,x_{r})\) is \(M_{p}\) -regular.

In what follows we shall consider mainly the notion of M-regular sequence in the case where the ring A is local Noetherian, the module M is finitely generated and the elements of the sequence belong to \(m_{A}\) ; the notion of M-regular sequence then coincides with that of completely secant sequence for M (A, X, p. 160, cor. 1).

PROPOSITION 7.--- Let A be a ring, J an ideal of A, M an A-module, and \((x_{1},\ldots,x_{r})\) an M-regular sequence of elements of J. One has

\[
\operatorname{prof} _ {\mathrm{A}} (\mathrm{J}; \mathrm{M}) = r + \operatorname{prof} _ {\mathrm{A}} (\mathrm{J}; \mathrm{M} / (x _ {1} \mathrm{M} + \dots + x _ {r} \mathrm{M}))
\]

and in particular \(\operatorname{prof}_{\mathrm{A}}(\mathrm{J};\mathrm{M})\geqslant r\)

The case \(r=1\) follows from remark 5 of n° 1, applied to the exact sequence

\[
0 \rightarrow \mathrm{M} \xrightarrow {(x _ {1}) _ {\mathrm{M}}} \mathrm{M} \longrightarrow \mathrm{M} / x _ {1} \mathrm{M} \rightarrow 0.
\]

The general case is deduced from it by induction on r.

THEOREM 2.--- Let A be a Noetherian ring, J an ideal of A, and M a finitely generated A-module.

\begin{enumerate}
\def\labelenumi{\alph{enumi})}
\item
  Suppose that \(\mathrm{prof}_{\mathrm{A}}(\mathrm{J};\mathrm{M})\) is finite. Then every M-regular sequence of elements of J can be extended to an M-regular sequence of length \(\mathrm{prof}_{\mathrm{A}}(\mathrm{J};\mathrm{M})\) of elements of J.
\item
  The depth of M with respect to J is the supremum of the lengths of M-regular sequences formed from elements of J.
\item
  For \(\operatorname{prof}_{\mathrm{A}}(\mathrm{J};\mathrm{M})\) to to be finite, it is necessary and sufficient that the support of M meet \(\mathrm{V}(\mathrm{J})\) , or equivalently that one has \(\mathrm{JM} \neq \mathrm{M}\) .
\end{enumerate}

Let \((x_{1},\ldots,x_{r})\) be an M-regular sequence of elements of J. One has \(r\leqslant\operatorname{prof}_{\mathrm{A}}(\mathrm{J};\mathrm{M})\) (prop. 7); suppose that the inequality is strict. Let N denote the A-module \(\mathrm{M}/(x_{1}\mathrm{M}+\ldots+x_{r}\mathrm{M})\) . One has \(\operatorname{prof}_{\mathrm{A}}(\mathrm{J};\mathrm{N})>0\) (loc. cit.), so that there exists an element x of J such that the homothety \(x_{N}\) is injective (n° 1, remark 2), that is, the sequence \((x_{1},\ldots,x_{r},x)\) is M-regular. It follows by induction that for every integer s such that \(r\leqslant s\leqslant\operatorname{prof}_{\mathrm{A}}(\mathrm{J};\mathrm{M})\) the sequence \((x_{1},\ldots,x_{r})\) can be completed to an M-regular sequence of length s, which entails assertions a) and b). Assertion c) follows from remark 1 of no. 1 and cor. 1 of th. 1 of no. 3.

COROLLARY 1. For every M-regular sequence \((x_{1},\ldots ,x_{r})\) of elements of J, the following properties are equivalent:

\begin{enumerate}
\def\labelenumi{(\roman{enumi})}
\item
  we have \(r = \mathrm{prof}_{\mathrm{A}}(\mathrm{J};\mathrm{M})\)
\item
  the sequence \((x_{1},\ldots,x_{r})\) is maximal among the M-regular sequences formed from elements of J;
\item
  the A-module \(\mathrm{M}/(x_{1}\mathrm{M}+\ldots+x_{r}\mathrm{M})\) has a nonzero element annihilated by J ;
\end{enumerate}

\[
\text {(iv)} \text {   we have   } \operatorname{Ass} \bigl (\mathrm{M} / (x _ {1} \mathrm{M} + \ldots + x _ {r} \mathrm{M}) \bigr) \cap \mathrm{V} (\mathrm{J}) \neq \emptyset .
\]

The equivalence of (i) and (ii) follows from Th. 2; the equivalence of (ii), (iii), and (iv) follows from Remark 2 of No. 1 applied to the A-module M/(x₁M + \ldots{} + xᵣM).

COROLLARY 2. --- Let A be a Noetherian local ring, and let M be a nonzero finitely generated A-module. Then

\[
\operatorname{prof} _ {\mathrm{A}} (\mathrm{M}) \leqslant \dim_ {\mathrm{A}} (\mathrm{M}) <   + \infty .
\]

Indeed every M-regular sequence of elements of \(\mathfrak{m}_{\mathrm{A}}\) is completely secant for M (A, X, p.~157, prop. 5), hence secant for M (VIII, § 3, n° 2, cor. of prop. 3); consequently its length is bounded above by the integer \(\dim_{\Lambda}(\mathrm{M})\) (loc. cit., th. 1).

\subsection{Depth along a closed subset}

Let A be a Noetherian ring, F a closed subset of \(\operatorname{Spec}(A)\) and M an A-module. By cor. 2 of prop. 2 of no. 1, the element \(\operatorname{prof}_{A}(J;M)\) of \(\mathbf{N} \cup \{+\infty\}\) does not depend on the ideal J of A such that \(F = V(J)\) ; it is called the depth of M along F and is denoted by \(\operatorname{prof}_{F}(M)\) .

Remarks.-- 1) Let \(r\) be an integer. By prop. 2 of no. 1 and II, § 4, no. 4, cor. 2 of prop. 17, the inequality \(\mathrm{prof}_{\mathrm{F}}(\mathbf{M}) \geqslant r\) is equivalent to the following property: for every finitely generated A-module N whose support is contained in F, we have \(\mathrm{Ext}_{\mathrm{A}}^{i}(\mathrm{N},\mathrm{M}) = 0\) for \(i < r\) .

\begin{enumerate}
\def\labelenumi{\arabic{enumi})}
\setcounter{enumi}{1}
\tightlist
\item
  Suppose that the A-module M is finitely generated. By Remark 2 of No. 1 and Th. 2 of No. 4, we have the following equivalences
\end{enumerate}

\[
\begin{array}{c} \operatorname{prof} _ {\mathrm{F}} (\mathrm{M}) = 0 \Longleftrightarrow \operatorname{Ass} (\mathrm{M}) \cap \mathrm{F} \neq \emptyset \\ \operatorname{prof} _ {\mathrm{F}} (\mathrm{M}) <   + \infty \Longleftrightarrow \operatorname{Supp} (\mathrm{M}) \cap \mathrm{F} \neq \emptyset . \end{array}
\]

PROPOSITION 8. Let A be a Noetherian ring, and let F be a closed subset of \(\operatorname{Spec}(\Lambda)\) and let M be a finitely generated A-module. We have

\[
\operatorname{prof} _ {F} (M) = \inf _ {\mathfrak {p} \in F} \operatorname{prof} _ {A _ {\mathfrak {p}}} \left(M _ {\mathfrak {p}}\right) = \inf _ {\mathfrak {p} \in \operatorname{Supp} (M) \cap F} \operatorname{prof} _ {A _ {\mathfrak {p}}} \left(M _ {\mathfrak {p}}\right).
\]

This is clear if \(\operatorname{prof}_{F}(M) = +\infty\) (remark 2). If \(\operatorname{prof}_{F}(M) = 0\) , there exists a prime ideal \(\mathfrak{p} \in \operatorname{Ass}(M) \cap F\) (remark 2); we have \(\mathfrak{p}A_{\mathfrak{p}} \in \operatorname{Ass}(M_{\mathfrak{p}})\) (IV, § 1, no. 2, prop. 5), hence \(\operatorname{prof}_{\Lambda_{\mathfrak{p}}}(\mathrm{M}_{\mathfrak{p}}) = 0\) (remark 2 of No. 1), whence the proposition in this case.

Suppose \(0 < \operatorname{prof}_{\mathrm{F}}(\mathrm{M}) < +\infty\). Let J be an ideal of A such that \(\mathrm{V}(\mathrm{J}) = \mathrm{F}\) , and let \(x\) an element of J such that the homothety \(x_{\mathrm{M}}\) is injective (loc. cit.). For each prime ideal \(\mathfrak{p}\) the homothety \(x_{\mathrm{M}_{\mathfrak{p}}}\) is injective. By prop. 7 of no. 4, we therefore have

\[
\begin{array}{c} \operatorname{prof} _ {\mathrm{F}} (\mathrm{M} / x \mathrm{M}) = \operatorname{prof} _ {\mathrm{F}} (\mathrm{M}) - 1 \\ \operatorname{prof} _ {\Lambda_ {\mathfrak {p}}} \big ((\mathrm{M} / x \mathrm{M}) _ {\mathfrak {p}} \big) = \operatorname{prof} _ {\Lambda_ {\mathfrak {p}}} (\mathrm{M} _ {\mathfrak {p}}) - 1. \end{array}
\]

We then conclude by induction on the integer \(\mathrm{prof}_{\mathbb{F}}(\mathbb{M})\) .

Remark 3. If q is a point of Supp(M), we therefore have \(\operatorname{prof}_{\mathrm{A}}(\mathfrak{q};\mathrm{M}) = \inf_{\mathfrak{p}\supseteq q} \operatorname{prof}_{\Lambda_{\mathfrak{p}}}(\mathrm{M}_{\mathfrak{p}})\) In particular, we have the inequality \(\operatorname{prof}_{\mathrm{A}}(\mathfrak{q};\mathrm{M}) \leqslant \operatorname{prof}_{\mathrm{A}_{\mathfrak{q}}}(\mathrm{M}_{\mathfrak{q}})\) ; equality holds when q is maximal. In the general case, one may have \(\operatorname{prof}_{\mathrm{A}}(\mathfrak{q};\mathrm{M}) < \operatorname{prof}_{\Lambda_{\mathfrak{q}}}(\mathrm{M}_{\mathfrak{q}})\) ; one may also have \(\operatorname{prof}_{\mathrm{A}}(\mathfrak{q};\mathrm{M}) < \inf \operatorname{prof}_{\mathrm{A}_{\mathfrak{m}}}(\mathrm{M}_{\mathfrak{m}})\) where m ranges over the set of maximal ideals of A containing q. Let, for example, p be a nonmaximal prime ideal of A, containing q and distinct from q; set M = A/p. One has \(\operatorname{prof}_{\mathrm{A}}(\mathfrak{q};\mathrm{M}) = 0\) , \(\operatorname{prof}_{\mathrm{A}_{\mathfrak{q}}}(\mathrm{M}_{\mathfrak{q}}) = +\infty\) and \(\operatorname{prof}_{\mathrm{A}_{\mathfrak{m}}}(\mathrm{M}_{\mathfrak{m}}) > 0\) for every maximal ideal m of A.

PROPOSITION 9.--- Let A be a Noetherian ring, M and N two finitely generated A-modules, and F the support of N. Then \(\operatorname{prof}_{\mathbf{F}}(\mathbf{M})\) is the lower bound (in \(\mathbf{N} \cup \{+\infty\}\) ) of the set of integers \(n\) such that \(\operatorname{Ext}_{\Lambda}^{n}(\mathbf{N}, \mathbf{M}) \neq 0\) .

By Remark 1, we have \(\mathrm{Ext}_{\mathrm{A}}^{i}(\mathrm{N},\mathrm{M}) = 0\) for every \(i < \mathrm{prof}_{\mathrm{F}}(\mathrm{M})\) . It remains to prove that if \(\mathrm{prof}_{\mathrm{F}}(\mathrm{M}) = n < +\infty\) , one has \(\mathrm{Ext}_{\mathrm{A}}^{n}(\mathrm{N},\mathrm{M}) \neq 0\) . Let J be the annihilator of N; we have \(\mathrm{F} = \mathrm{V}(\mathrm{J})\) , hence \(\mathrm{prof}_{\mathrm{F}}(\mathrm{M}) = \mathrm{prof}_{\mathrm{A}}(\mathrm{J};\mathrm{M})\) . By cor. 1 of th. 2 ( \(n^{\circ}4\) ), there exists an M-regular sequence ( \(x_{1},\ldots,x_{n}\) ) of length \(n\) consisting of elements of J, and the depth of the A-module relative to J \(\overline{\mathrm{M}} = \mathrm{M}/(x_{1}\mathrm{M}+\ldots+x_{n}\mathrm{M})\) is zero. By A, X, p. 166, prop. 9, it suffices to prove that \(\mathrm{Hom}_{\mathrm{A}}(\mathrm{N},\overline{\mathrm{M}})\) is nonzero. Now, by prop. 8, there exists \(\mathfrak{p} \in \mathrm{Supp}(\mathrm{M}) \cap \mathrm{Supp}(\mathrm{N})\) such that \(\mathrm{prof}_{\mathrm{A}_{\mathfrak{p}}}(\overline{\mathrm{M}}_{\mathfrak{p}}) = 0\) , that is \(\mathrm{Hom}_{\mathrm{A}_{\mathfrak{p}}}(\kappa(\mathfrak{p}),\overline{\mathrm{M}}_{\mathfrak{p}}) \neq 0\) . Since \(\mathbf{N}_{\mathfrak{p}}\) is nonzero, the \(\kappa(\mathfrak{p})\) -vector space \(\mathbf{N}_{\mathfrak{p}} \otimes_{\mathbf{A}_{\mathfrak{p}}} \kappa(\mathfrak{p})\) is nonzero (Nakayama's lemma), as is its dual; there therefore exists a surjective \(\mathbf{A}_{\mathfrak{p}}\) -linear map from \(\mathbf{N}_{\mathfrak{p}}\) in \(\kappa(\mathfrak{p})\) . It follows that one has \(\mathrm{Hom}_{\mathrm{A}_{\mathfrak{p}}}(\mathbf{N}_{\mathfrak{p}},\overline{\mathbf{M}}_{\mathfrak{p}}) \neq 0\) , hence \(\mathrm{Hom}_{\mathrm{A}}(\mathrm{N},\overline{\mathrm{M}}) \neq 0\) (II, § 2, \(n^{\circ}7\) , prop. 19), which is what we wanted to prove.

Remark 4.--- Let A be a Noetherian ring, N a finitely generated A-module. One sometimes calls the grade of N, and denotes it by grade (N), the greatest lower bound in \(\mathbf{N}\cup\{+\infty\}\) of the set of integers \(n\) such that \(\operatorname{Ext}_{\mathrm{A}}^{n}(\mathrm{N},\mathrm{A})\) is nonzero. By prop. 9, this is also the depth of A along the support of N, or again ( \(n^{\circ}\) 4, th. 2) the least upper bound of the set of lengths of A-regular sequences of elements of the annihilator of N. Since for every prime ideal p of A the annihilator of \(N_{p}\) is equal to \(\operatorname{Ann}(N)_{p}\) (II, § 2, \(n^{\circ}\) 4, formula (9)), one deduces from prop. 5 of \(n^{\circ}\) 3 the equality

\[
\operatorname{grade} (\mathrm{N}) = \inf _ {\mathfrak {p} \in \operatorname{Spec} (\Lambda)} \operatorname{grade} \left(\mathrm{N} _ {\mathfrak {p}}\right) = \inf _ {\mathfrak {m} \in \Omega} \operatorname{grade} \left(\mathrm{N} _ {\mathfrak {m}}\right),
\]

where Ω denotes the set of maximal ideals of A.

\subsection{Depth of algebras}

Lemma 2.--- Let \(\rho: A \to B\) be a local homomorphism of Noetherian local rings, N a finitely generated B-module and y an element of \(\mathfrak{m}_{\mathrm{B}}\) . The following two conditions are equivalent:

\begin{enumerate}
\def\labelenumi{(\roman{enumi})}
\item
  the A-module N/yN is flat and the homothety \(y_{N}\) is injective;
\item
  the A-module N is flat and the homothety \(y_{\kappa_{A}\otimes N}\) is injective.
\end{enumerate}

When they are satisfied, the homothety \(y_{\mathrm{M} \otimes_{\Lambda}\mathrm{N}}\) is injective for every A-module M.

Suppose the hypotheses of (i) are satisfied, and prove (ii) as well as the last assertion. Let M be an A-module. Since the A-module N/yN is flat, one deduces from the exact sequence \(0 \rightarrow N \xrightarrow{y_{N}} N \longrightarrow N/yN \rightarrow 0\) the exact sequences

\[
0 \rightarrow \mathrm{M} \otimes_ {\mathrm{A}} \mathrm{N} \xrightarrow {u} \mathrm{M} \otimes_ {\mathrm{A}} \mathrm{N} \longrightarrow \mathrm{M} \otimes_ {\mathrm{A}} (\mathrm{N} / y \mathrm{N}) \rightarrow 0
\]

\[
0 \to \operatorname{Tor} _ {1} ^ {\mathrm{A}} (\mathrm{M}, \mathrm{N}) \stackrel {v} {\longrightarrow} \operatorname{Tor} _ {1} ^ {\mathrm{A}} (\mathrm{M}, \mathrm{N}) \to 0
\]

where \(u = 1_{\mathrm{M}} \otimes y_{\mathrm{N}}\) and \(v = \operatorname{Tor}_1^{\mathrm{A}}(1_{\mathrm{M}}, y_{\mathrm{N}})\) ; it follows that the homothety of ratio \(y\) is injective in \(\mathrm{M} \otimes_{\mathrm{A}} \mathrm{N}\) , and bijective in \(\operatorname{Tor}_1^{\mathrm{A}}(\mathrm{M}, \mathrm{N})\) . Suppose moreover that the A-module \(\mathrm{M}\) is of finite type; then the B-module \(\operatorname{Tor}_1^{\mathrm{A}}(\mathrm{M}, \mathrm{N})\) is of finite type (A, X, p. 107, prop. 6), hence is zero since \(y\) belongs to \(\mathfrak{m}_{\mathrm{B}}\) (Nakayama's lemma), which implies that the A-module \(\mathrm{N}\) is flat (A, X, p. 74, th. 2).

(ii)⇒(i): suppose the hypotheses of (ii) are satisfied. Consider the two exact sequences of B-modules

\begin{enumerate}
\def\labelenumi{(\arabic{enumi})}
\tightlist
\item
\end{enumerate}

\[
0 \to \mathrm{Ker} (y _ {\mathrm{N}}) \longrightarrow \mathrm{N} \stackrel {p} {\longrightarrow} \mathrm{Im} (y _ {\mathrm{N}}) \to 0\tag{1}
\]

\[
0 \to \mathrm{Im} (y _ {\mathrm{N}}) \stackrel {i} {\longrightarrow} \mathrm{N} \longrightarrow \mathrm{N} / y \mathrm{N} \to 0,\tag{2}
\]

where \(p\) and \(i\) are the canonical homomorphisms. It follows that the homomorphism \(1 \otimes p : \kappa_{\mathrm{A}} \otimes_{\Lambda} \mathrm{N} \longrightarrow \kappa_{\mathrm{A}} \otimes_{\mathrm{A}} \mathrm{Im}(y_{\mathrm{N}})\) is surjective, and (since N is flat) that the kernel of the homomorphism \(1 \otimes i : \kappa_{\mathrm{A}} \otimes_{\mathrm{A}} \mathrm{Im}(y_{\mathrm{N}}) \longrightarrow \kappa_{\mathrm{A}} \otimes_{\Lambda} \mathrm{N}\) is isomorphic to \(\operatorname{Tor}_1^{\mathrm{A}}(\kappa_{\mathrm{A}}, \mathrm{N}/y\mathrm{N})\) . But the map \((1 \otimes i) \circ (1 \otimes p)\) , equal to \(y_{\kappa_{\mathrm{A}} \otimes_{\mathrm{A}} \mathrm{N}}\) , is injective by hypothesis; it follows that \(1 \otimes p\) is bijective and \(1 \otimes i\) injective, and consequently one has \(\operatorname{Tor}_1^{\mathrm{A}}(\kappa_{\mathrm{A}}, \mathrm{N}/y\mathrm{N}) = 0\) . It follows that the A-module \(\mathrm{N}/y\mathrm{N}\) is flat (III, § 5, no. 2, th. 1 and no. 4, prop. 2).

Since N and N/yN are flat over A, the same holds for \(\mathrm{Im}(y_{\mathrm{N}})\) (exact sequence (2)). We then deduce from the exact sequence (1) that \(\kappa_{A} \otimes_{A} \mathrm{Ker}(y_{\mathrm{N}})\) is isomorphic to the kernel of \(1 \otimes p\) ; hence it is zero; thus the homothety \(y_{N}\) is injective by Nakayama's lemma.

PROPOSITION 10. Let \(\rho: \Lambda \to B\) a local homomorphism of Noetherian local rings, N a finitely generated B-module, and \(y = (y_1, \ldots, y_s)\) be a sequence of elements of \(\mathfrak{m}_{\mathrm{B}}\) . Let \(\mathfrak{y}\) the ideal of B generated by that sequence. The following conditions are equivalent:

\begin{enumerate}
\def\labelenumi{(\roman{enumi})}
\item
  the A-module N/\(\mathfrak{y}\)N is flat and the sequence \(y\) is N-regular;
\item
  the A-module N is flat and the sequence \(y\) is \((\kappa_{A} \otimes_{A} N)\)-regular.
\end{enumerate}

When they are satisfied, for every A-module M, the sequence \(y\) is \(M\otimes_{\mathrm{A}}N\)-regular.

Let us prove the equivalence of (i) and (ii) by induction on s. The case s = 0 being evident, suppose \(s \geqslant 1\) . Let us denote by \(y'\) the sequence \((y_{1}, \ldots, y_{s-1})\) and \(\mathfrak{y}'\) the ideal of B that it generates. By Lemma 2 applied to the B-module N/ \(y'\) N and to the element \(y_{s}\) of B, condition (i) is equivalent to

(i') the A-module N/ \(\mathfrak{y}'\) N is flat, and the sequence \(y'\) is N-regular and the homothety with ratio \(y_s\) in \(\kappa_{\mathrm{A}} \otimes_{\mathrm{A}} (\mathrm{N}/\mathfrak{y}'\mathrm{N}) = (\kappa_{\mathrm{A}} \otimes_{\mathrm{A}} \mathrm{N})/\mathfrak{y}'(\kappa_{\mathrm{A}} \otimes_{\mathrm{A}} \mathrm{N})\) is injective.

This condition is equivalent to (ii) by the induction hypothesis.

The last assertion follows similarly by induction on \(s\) from the last assertion of Lemma 2.

PROPOSITION 11. Let \(\rho: A \to B\) be a local homomorphism of Noetherian local rings, M a finitely generated A-module and N a finitely generated B-module; we assume that the A-module N is flat.

\begin{enumerate}
\def\labelenumi{\alph{enumi})}
\item
  Let \((x_{1},\ldots,x_{r})\) be a sequence of elements of \(m_{A}\) regular for the A-module M, and \((y_{1},\ldots,y_{s})\) be a sequence of elements of \(m_{B}\) regular for the B-module \(\kappa_{A}\otimes_{A}N\) ; then the sequence \((y_{1},\ldots,y_{s},\rho(x_{1}),\ldots,\rho(x_{r}))\) of elements of \(m_{B}\) is regular for the B-module \(M\otimes_{A}N\) .
\item
  We have the equality
\end{enumerate}

\[
\mathrm{prof} _ {\mathrm{B}} (\mathrm{M} \otimes_ {\mathrm{A}} \mathrm{N}) = \mathrm{prof} _ {\mathrm{A}} (\mathrm{M}) + \mathrm{prof} _ {\mathrm{B}} (\kappa_ {\mathrm{A}} \otimes_ {\mathrm{A}} \mathrm{N}).
\]

Let \(\mathfrak{x}\) be the ideal of A generated by the sequence \(\mathbf{x}\) and \(\mathfrak{y}\) be the ideal of B generated by \(y\) . By prop. 10, the sequence \(y\) is \(M \otimes_A N\) -regular and \(N/\mathfrak{y}N\) is flat over \(A\) , so that the sequence \(\boldsymbol{\rho}(\mathbf{x}) = (\boldsymbol{\rho}(x_1), \ldots, \boldsymbol{\rho}(x_r))\) is regular for \(M \otimes_A (N/\mathfrak{y}N) = (M \otimes_A N)/\mathfrak{y}(M \otimes_A N)\) . This proves a).

To prove b), we may assume M and N are nonzero. By Nakayama's lemma, \(\kappa_{\mathrm{A}} \otimes_{\mathrm{A}} \mathrm{N}\) is also nonzero, so that \(\operatorname{prof}_{\mathrm{A}}(\mathrm{M})\) and \(\operatorname{prof}_{\mathrm{B}}(\kappa_{\mathrm{A}} \otimes_{\mathrm{A}} \mathrm{N})\) are finite (no. 4, cor. 2 of th. 2). Take the maximal regular sequences x and y; we then have \(r = \operatorname{prof}_{\mathrm{A}}(\mathrm{M})\) , \(s = \operatorname{prof}_{\mathrm{B}}(\kappa_{\mathrm{A}} \otimes_{\mathrm{A}} \mathrm{N})\) and there exists an injective A-linear map \(u: \kappa_{\mathrm{A}} \to \mathrm{M} / x \mathrm{M}\) and an injective B-linear map \(v: \kappa_{\mathrm{B}} \to \kappa_{\mathrm{A}} \otimes_{\mathrm{A}} (\mathrm{N} / \mathfrak{y} \mathrm{N})\) . Since N/ \(\mathfrak{y}\) N is flat over A, the B-linear map \((u \otimes 1_{\mathrm{N} / \mathfrak{y} \mathrm{N}}) \circ v\) of \(\kappa_{\mathrm{B}}\) into (M/xM) \(\otimes_{\mathrm{A}} (\mathrm{N} / \mathfrak{y} \mathrm{N}) = (\mathrm{M} \otimes_{\mathrm{A}} \mathrm{N}) / (\rho(x) + \mathfrak{y})(\mathrm{M} \otimes_{\mathrm{A}} \mathrm{N})\) is injective. This implies the equality \(\operatorname{prof}_{\mathrm{B}}(\mathrm{M} \otimes_{\mathrm{A}} \mathrm{N}) = r + s\) (loc. cit.).

Remark. -- Recall that, under the preceding hypotheses, we have

\[
\dim_ {B} (M \otimes_ {A} N) = \dim_ {A} (M) + \dim_ {B} (\kappa_ {A} \otimes_ {A} N)
\]

(VIII, § 3, no. 4, prop. 7).

COROLLARY.--- Let \(\rho: A \to B\) a local homomorphism of Noetherian local rings making B a flat A-module. We have

\[
\begin{array}{l} \operatorname{prof} (\mathrm{B}) = \operatorname{prof} (\mathrm{A}) + \operatorname{prof} \left(\kappa_ {\Lambda} \otimes_ {\Lambda} \mathrm{B}\right), \\ \dim (\mathrm{B}) = \dim (\mathrm{A}) + \dim \left(\kappa_ {\mathrm{A}} \otimes_ {\mathrm{A}} \mathrm{B}\right). \end{array}
\]

Indeed, the depth (resp. the dimension) of the B-module \(\kappa_{\mathrm{A}}\otimes_{\mathrm{A}}\mathrm{B}\) is equal to the depth (resp. the dimension) of the ring \(\kappa_{\mathrm{A}}\otimes_{\mathrm{A}}\mathrm{B}\) by the corollary of Prop. 4 (resp. by VIII, § 1, no. 4).

\subsection{Bounds on depth}

PROPOSITION 12.--- Let A be a Noetherian local ring, M a nonzero finitely generated A-module, and J an ideal of A distinct from A. We have the sequence of inequalities

\[
\begin{array}{r l} \operatorname{prof} _ {A} (J; M) & \leqslant \operatorname{codim} (\operatorname{Supp} (M) \cap V (J), \operatorname{Supp} (M)) \leqslant \dim (M) - \dim (M / J M) \\ & \leqslant [ J / \mathfrak {m} _ {A} J: \kappa_ {A} ]. \end{array}
\]

For every \(\mathfrak{p} \in \operatorname{Supp}(\mathrm{M})\cap \mathrm{V}(\mathrm{J})\), \(\operatorname{prof}_{\mathrm{A}}(\mathrm{J};\mathrm{M})\) is less than \(\dim_{\mathrm{A}_{\mathfrak{p}}}(\mathrm{M}_{\mathfrak{p}})\) (No. 5, Prop. 8 and No. 4, Cor. 2 of Th. 2), that is, (VIII, § 1, No. 4, Prop. 9)

to \(\operatorname{codim}(V(\mathfrak{p}),\operatorname{Supp}(\mathrm{M}))\). When \(\mathfrak{p}\) ranges over \(\operatorname{Supp}(\mathrm{M}) \cap \mathrm{V}(\mathrm{J})\), \(V(\mathfrak{p})\) describes the irreducible closed subsets of \(\operatorname{Supp}(\mathrm{M}) \cap \mathrm{V}(\mathrm{J})\), whence the first inequality. The second follows from Prop. 3 of VIII, § 1, No. 2. Moreover, one can find a generating set of J of cardinality \([J/\mathfrak{m}_{\mathrm{A}}J:\kappa_{\mathrm{A}}]\) (II, § 3, No. 2, Cor. 2 of Prop. 4); the third inequality then follows from VIII, § 3, No. 2, formula (8).

Remark 1.-- Consider the chain of inequalities in Prop. 12.

\begin{enumerate}
\def\labelenumi{\alph{enumi})}
\item
  For us to have \(\mathrm{prof}_{\mathrm{A}}(\mathrm{J};\mathrm{M}) = [\mathrm{J} / \mathfrak{m}_{\mathrm{A}}\mathrm{J}:\kappa_{\mathrm{A}}]\) it is necessary and sufficient that J can be generated by an M-regular sequence (no. 3, cor. 2 of th. 1 and A, X, p. 160, cor. 1).
\item
  The equality \(\dim(\mathrm{M}) - \dim(\mathrm{M}/\mathrm{JM}) = [\mathrm{J}/\mathfrak{m}_{\Lambda}\mathrm{J} : \kappa_{\Lambda}]\) means that J can be generated by a regular sequence for M (VIII, § 3, no. 2).
\item
  If M is Macaulay, we have \(\operatorname{prof}_{\mathrm{A}}(\mathbf{J};\mathbf{M}) = \dim (\mathbf{M}) - \dim (\mathbf{M} / \mathbf{JM})\) (§ 2, no. 2, cor. of prop. 2).
\end{enumerate}

Lemma 3.--- Let A be a Noetherian ring, \(\mathfrak{p} \subset \mathfrak{p}_1 \subset \ldots \subset \mathfrak{p}_{r-1} \subset \mathfrak{q}\) a saturated chain of length \(r\) of prime ideals of A, 'M a \(\Lambda\) -module of finite type and \(n\) an integer. If the A-module \(\mathrm{Ext}_{\mathrm{A}_{\mathfrak{p}}}^{n}(\kappa(\mathfrak{p}), \mathrm{M}_{\mathfrak{p}})\) is nonzero, the same holds for \(\mathrm{Ext}_{\mathrm{A}_{\mathfrak{q}}}^{n+r}(\kappa(\mathfrak{q}), \mathrm{M}_{\mathfrak{q}})\) .

It obviously suffices to treat the case \(r = 1\); substituting then A, M, p, and q by \(A_q\) , \(M_q\) , \(pA_q\) and \(qA_q\) respectively, we are reduced to treating the case where A is local and where \(q = m_A\) . Let \(x\) be an element of \(m_A - p\) . The \(A_p\) -module \(\text{Ext}_A^n(A/p, M) \otimes_A A_p\) is isomorphic to \(\text{Ext}_{A_p}^n(\kappa(p), M_p)\) (A, X, p. 111, prop. 10 b)), hence is nonzero by hypothesis; a fortiori \(\text{Ext}_A^n(A/p, M)\) is not zero. The exact sequence

\[
0 \rightarrow \mathrm{A} / \mathfrak {p} \stackrel {{x _ {\mathrm{A} / \mathfrak {p}}}} {{\longrightarrow}} \mathrm{A} / \mathfrak {p} \longrightarrow \mathrm{A} / (\mathfrak {p} + x \mathrm{A}) \rightarrow 0
\]

provides an exact sequence of extension modules

\[
\operatorname{Ext} _ {\mathrm{A}} ^ {n} (\mathrm{A} / \mathfrak {p}, \mathrm{M}) \xrightarrow {u} \operatorname{Ext} _ {\mathrm{A}} ^ {n} (\mathrm{A} / \mathfrak {p}, \mathrm{M}) \longrightarrow \operatorname{Ext} _ {\mathrm{A}} ^ {n + 1} (\mathrm{A} / (\mathfrak {p} + x \mathrm{A}), \mathrm{M}),
\]

where \(u\) is the homothety of ratio \(x\) . By Nakayama's lemma, this is not surjective, so the A-module \(\text{Ext}_A^{n+1}(\text{A}/(\mathfrak{p} + x\text{A}), \text{M})\) is not zero. But the only prime ideal of A containing \(\mathfrak{p} + x\text{A}\) is \(\mathfrak{m}_\text{A}\); therefore the A-module \(\text{A}/(\mathfrak{p} + x\text{A})\) has finite length (VIII, § 3, no. 2, lemma 2). If \(\text{Ext}_\text{A}^{n+1}(\kappa_\text{A}, \text{M})\) were zero, one would deduce, by induction on the length of N, the vanishing of \(\text{Ext}_\text{A}^{n+1}(\text{N}, \text{M})\) for every A-module N of finite length. This contradiction completes the proof.

PROPOSITION 13.--- Let A be a Noetherian ring, M a finitely generated A-module, and p and q two prime ideals in Supp(M) with p ⊂ q. Then

\[
\operatorname{prof} _ {A _ {\mathfrak {q}}} (M _ {\mathfrak {q}}) \leqslant \operatorname{prof} _ {A _ {\mathfrak {p}}} (M _ {\mathfrak {p}}) + \dim (A _ {\mathfrak {q}} / \mathfrak {p} A _ {\mathfrak {q}}).
\]

More precisely, for every saturated chain of prime ideals \(\mathfrak{p} \subset \mathfrak{p}_1 \subset \ldots \subset \mathfrak{p}_{r-1} \subset \mathfrak{q}\) , one has \(\operatorname{prof}_{\mathrm{A}_{\mathfrak{q}}}(\mathrm{M}_{\mathfrak{q}}) \leqslant \operatorname{prof}_{\mathrm{A}_{\mathfrak{p}}}(\mathrm{M}_{\mathfrak{p}}) + r\) .

Set \(p = \text{prof}_{\text{A}_p}(\text{M}_p)\) and let us prove the second inequality. It is evident if \(p = +\infty\) otherwise, we have \(\text{Ext}_{\text{A}_p}^p(\kappa(\mathfrak{p}), \text{M}_p) \neq 0\) , whence \(\operatorname{Ext}_{\Lambda_{\mathfrak{q}}}^{p+r}(\kappa(\mathfrak{q}),\mathrm{M}_{\mathfrak{q}})\neq0\) by lemma 3, which implies \(\operatorname{prof}_{\mathrm{A}_{\mathfrak{q}}}(M_{\mathfrak{q}})\leqslant p+r\) . Since \(\dim(A_{\mathfrak{q}}/\mathfrak{p}A_{\mathfrak{q}})\) is the supremum of the lengths of saturated chains of prime ideals with given endpoints. \(\mathfrak{p}\) and \(\mathfrak{q}\) , the first assertion is a consequence of the second.

COROLLARY 1.--- We have the inequality

\[
\dim (M _ {\mathfrak {q}}) - \operatorname{prof} _ {\Lambda_ {\mathfrak {q}}} (M _ {\mathfrak {q}}) \geqslant \dim (M _ {\mathfrak {p}}) - \operatorname{prof} _ {A _ {\mathfrak {p}}} (M _ {\mathfrak {p}}) \geqslant 0.
\]

This follows from prop. 13 and the inequality \(\dim(M_{\mathfrak{q}}) \geqslant \dim(M_{\mathfrak{p}}) + \dim(A_{\mathfrak{q}} / \mathfrak{p}A_{\mathfrak{q}})\) (VIII, § 1, no. 4, prop. 9, a) and no. 3, prop. 7, b)).

COROLLARY 2.--- Let A be a Noetherian local ring and M a finitely generated A-module. Then one has the inequality

\[
\operatorname{prof} _ {\mathrm{A}} (\mathrm{M}) \leqslant \inf _ {\mathfrak {p} \in \operatorname{Ass} (\mathrm{M})} \dim (\mathrm{A} / \mathfrak {p}).
\]

Let \(\mathfrak{p}\) a prime ideal associated with M; we have \(\mathrm{prof}_{\mathrm{A}_{\mathfrak{p}}}(\mathrm{M}_{\mathfrak{p}})=0\) (No. 1, Remark 2). Prop. 13 applied to the ideals \(\mathfrak{p}\subset\mathfrak{m}_{\mathrm{A}}\) provides the inequality \(\mathrm{prof}_{\mathrm{A}}(\mathrm{M})\leqslant\dim(\mathrm{A}/\mathfrak{p})\) , whence the corollary.

Remark 2.--- We have \(\sup_{\mathfrak{p}\in\operatorname{Ass}(\mathbb{M})}\dim(\mathrm{A}/\mathfrak{p})=\dim(\mathrm{M})\) (VIII, § 1, no. 4, remark 2) and one recovers the inequality \(\operatorname{prof}(\mathrm{M})\leqslant\dim(\mathrm{M})\) for \(\mathrm{M}\neq0\) (No. 4, cor. 2 of th. 2).

\subsection{Locally integral Noetherian rings; normal Noetherian rings}

Let A be a Noetherian ring. We denote by \((\mathrm{Y}_{i})_{i\in\mathrm{I}}\) the finite family (II, § 4, no. 2, prop. 10 and no. 3, cor. 7 of prop. 11) of connected components of Spec(A). According to II, § 4, no. 3, prop. 15, for each i there exists a unique idempotent element \(e_{i}\) of A such that \(\mathrm{Y}_{i}=\mathrm{V}(e_{i})\) , and the canonical map from A to the product of the rings A/Ae \(_{i}\) is bijective. The quotient rings A/Ae \(_{i}\) of A are said to be the canonical components of A. Set \(f_{i}=1-e_{i}\) . One has \(\sum_{i\in I}f_{i}=1\) and \((f_{i})_{i\in I}\) is an orthogonal family of nonzero idempotents of A (loc. cit.). It follows that the image of \(f_{i}\) in A/Ae \(_{j}\) is equal to 1 if j=i and to 0 if j≠i; one therefore deduces from the ring homomorphism A→ \(\prod_{j}A/Ae_{j}\) a canonical isomorphism A \(_{f_{i}}\rightarrow A/Ae_{i}\) .

According to loc. cit., cor. 2 of prop. 14, the following conditions are equivalent: (i) the connected components of Spec(A) are irreducible;

\begin{enumerate}
\def\labelenumi{(\roman{enumi})}
\setcounter{enumi}{1}
\item
  every prime (resp. maximal) ideal of A belongs to only one irreducible component of Spec(A);
\item
  every prime (resp. maximal) ideal of A contains only one minimal prime ideal;
\item
  for every prime (resp. maximal) ideal p of A, the topological space \(\operatorname{Spec}(A_{\mathfrak{p}})\) is irreducible;
\item
  for every canonical component C of A, the topological space \(\operatorname{Spec}(\mathbb{C})\) is irreducible.
\end{enumerate}

Let us now note that if A is reduced, all the rings \(A_{p}\) are reduced (II, § 2, n° 6, prop. 17) and that conversely, if the ring \(A_{m}\) is reduced for every maximal ideal m of A, then A is reduced (II, § 3, n° 3, cor. 2 of th. 1). Applying then II, § 4, n° 3, cor. 1 of prop. 14, one deduces from the preceding the equivalence of the following conditions:

\begin{enumerate}
\def\labelenumi{(\roman{enumi})}
\item
  A is reduced and the connected components of Spec(A) are irreducible;
\item
  for every prime (resp. maximal) ideal p of A, the ring A \(_{p}\) is an integral domain;
\item
  the canonical components of A are integral domains.
\end{enumerate}

A noetherian ring is said to be locally integral if it satisfies the equivalent conditions (i) to (iii) above.

Suppose the ring A is locally integral; let u be an isomorphism of A onto a (finite) product \(\prod_{j\in J}A_{j}\) of integral rings. There exists a bijection \(\sigma:J\to I\) such that the map from \(\operatorname{Spec}(\prod_{j\in J}A_{j})\) in \(\operatorname{Spec}(A)\) associated with u defines a homeomorphism of \(\operatorname{Spec}(A_{j})\) on the connected component \(Y_{\sigma(j)}\) of \(\operatorname{Spec}(A)\) ; we then deduce from u an isomorphism of the canonical component \(A/Ae_{\sigma(j)}\) on \(A_{j}\) .

PROPOSITION 14. - Let A be a Noetherian ring. The following conditions are equivalent:

\begin{enumerate}
\def\labelenumi{(\roman{enumi})}
\item
  A is reduced and integrally closed in its total ring of fractions;
\item
  A is isomorphic to the product of a finite family of integrally closed rings.
\item
  the canonical components of A are integrally closed;
\item
  for every prime (resp. maximal) ideal p of A, the ring \(A_{p}\) is integrally closed.
\end{enumerate}

The equivalence of (i) and (ii) follows from V, § 1, no. 2, cor. 2 of prop. 9, and that of (ii) and (iii) from the remarks preceding the proposition. Let p be a prime ideal of A; there exists a unique canonical component A' of A such that p belongs to the closed subset Spec(A') of Spec(A), and we have a canonical isomorphism \(A_{p} \simeq A'_{p}\). The equivalence of (iii) and (iv) therefore follows from loc. cit., no. 5, cor. 1 and cor. 3 of prop. 16.

A ring A is said to be normal if it is Noetherian and satisfies the equivalent conditions (i) to (iv) of Proposition 14. A Noetherian ring is integrally closed if and only if it is an integral domain and normal. A normal local ring is integrally closed.

\subsection{Depth and connectivity}

Lemma 4.--- Let A be a Noetherian ring, F a closed subset of Spec(A), U the complementary open set, and u : M → N a homomorphism of finitely generated A-modules.

\begin{enumerate}
\def\labelenumi{\alph{enumi})}
\item
  Suppose that \(u_{\mathfrak{p}}: \mathrm{M}_{\mathfrak{p}} \to \mathrm{N}_{\mathfrak{p}}\) is injective for all \(\mathfrak{p} \in \mathrm{U}\) and suppose we have \(\operatorname{prof}_{\mathrm{F}}(\mathrm{M}) \geqslant 1\). Then \(u\) is injective.
\item
  Suppose that \(u_{\mathfrak{p}}: \mathrm{M}_{\mathfrak{p}} \to \mathrm{N}_{\mathfrak{p}}\) is bijective for all \(\mathfrak{p} \in \mathrm{U}\) and suppose we have \(\operatorname{prof}_{\mathrm{F}}(\mathrm{M}) \geqslant 2\) and \(\operatorname{prof}_{\mathrm{F}}(\mathrm{N}) \geqslant 1\). Then u is bijective.
\item
  The hypotheses imply Supp(Ker \(u\) ) \(\subset\) F, then \(\operatorname{Hom}_{\Lambda}(\operatorname{Ker} u, \mathbf{M}) = 0\) ( \(n^{\circ}\) 5, remark 1); we therefore have \(\operatorname{Ker} u = 0\) .
\item
  We already know that \(u\) is injective, and we have \(\operatorname{Supp}(\operatorname{Coker} u) \subset \mathbb{F}\) . By loc. cit., we have \(\operatorname{Hom}_{\mathbf{A}}(\operatorname{Coker} u, \mathbf{N}) = 0\) and \(\operatorname{Ext}_{\mathbf{A}}^{1}(\operatorname{Coker} u, \mathbf{M}) = 0\) . From the exact sequence of extension modules
\end{enumerate}

\[
\operatorname{Hom} _ {\mathrm{A}} (\operatorname{Coker} u, \mathrm{N}) \rightarrow \operatorname{Hom} _ {\mathrm{A}} (\operatorname{Coker} u, \operatorname{Coker} u) \rightarrow \operatorname{Ext} _ {\mathrm{A}} ^ {1} (\operatorname{Coker} u, \mathrm{M})
\]

we deduce \(\operatorname{Hom}_{\mathrm{A}}(\operatorname{Coker} u, \operatorname{Coker} u) = 0\) , which implies \(\operatorname{Coker} u = 0\) .

Remark 1.—Let A be a Noetherian ring, F a closed subset of \(\operatorname{Spec}(A)\) , U the complementary open set. For one to have \(\operatorname{prof}_{\mathrm{F}}(A) \geqslant 1\) , it is necessary and sufficient that one has \(\operatorname{Ass}(A) \subset U\) (remark 2, no. 5). When this condition is satisfied, each irreducible component of \(\operatorname{Spec}(A)\) meets U, so that U is dense in \(\operatorname{Spec}(A)\) .

THEOREM 3 (Hartshorne).--- Let A be a Noetherian ring, F a closed subset of Spec(A), and U the complementary open set. Suppose that one has prof \(_{F}\) (A) ≥ 2. Then, for every connected component Y of Spec(A), the set Y ∩ U is connected and dense in Y.

Suppose first that \(\operatorname{Spec}(A)\) is connected. By Remark 1, \(U\) is dense in \(\operatorname{Spec}(A)\), and it remains to prove that it is connected. Suppose, for a contradiction, that two disjoint nonempty open subsets \(U_0\) and \(U_1\) of \(\operatorname{Spec}(A)\) have union \(U\). Since \(\operatorname{Ass}(A)\) is contained in \(U\) by Remark 1, it is a disjoint union of \(\operatorname{Ass}(A) \cap U_0\) and \(\operatorname{Ass}(A) \cap U_1\). By IV, § 1, No. 1, Prop. 4, there exist ideals \(J_0\) and \(J_1\) of A such that \(\operatorname{Ass}(J_i) = \operatorname{Ass}(A) \cap U_i\), \(\operatorname{Ass}(A/J_i) = \operatorname{Ass}(A) \cap U_{1-i}\quad (i=0,1)\). The complement of \(U_i\) in \(\operatorname{Spec}(A)\) contains \(\operatorname{Ass}(A/J_i)\) and \(\operatorname{Ass}(J_{1-i})\); since it is closed, it contains \(\operatorname{Supp}(A/J_i)\) and \(\operatorname{Supp}(J_{1-i})\). For \(\mathfrak{p} \in U_i\), we thus have \((J_i)_{\mathfrak{p}} = A_{\mathfrak{p}}\) and \((J_{1-i})_{\mathfrak{p}} = 0\); this implies in particular that \(J_0\) and \(J_1\) are distinct from A. Let B be the A-module \(A/J_0 \times A/J_1\) and let \(u: A \to B\) be the canonical homomorphism. By what precedes, the homomorphism \(u_{\mathfrak{p}}\) is bijective for every \(\mathfrak{p} \in U\); moreover, we have \(\operatorname{Ass}(B) \subset U_0 \cup U_1 = U\), hence \(\operatorname{prof}_F(B) \geqslant 1\) by Remark 1. Lemma 4 then implies that \(u\) is bijective, which contradicts the connectedness of \(\operatorname{Spec}(A)\).

Let us treat the general case. Let J be an ideal of A such that \(F = V(J)\), and let Y be a connected component of \(\operatorname{Spec}(A)\). By II, § 4, No. 3, Prop. 15, there exists an idempotent element \(f\) of A such that Y is identified with the part \(\operatorname{Spec}(A_f)\) of \(\operatorname{Spec}(A)\). Then \(Y \cap F\) is identified with \(V(J_f)\); we have \(\operatorname{prof}_{A_f}(J_f, A_f) \geqslant \operatorname{prof}_A(J; A) \geqslant 2\) by Prop. 6, a) of No. 3. It follows from the first part of the proof that \(Y \cap U = Y - (Y \cap F)\) is connected and dense in Y.

COROLLARY 1.--- The map which associates to each connected component of U its closure in Spec(A) is a bijection from the set of connected components of U onto the set of connected components of Spec(A).

COROLLARY 2.--- For every Noetherian local ring B of depth \(\geqslant\) 2, the topological space \(\operatorname{Spec}(\mathbf{B}) = \{\mathfrak{m}_{\mathbf{B}}\}\) is connected.

COROLLARY 3.--- Under the hypotheses of Theorem 3, suppose that \(\text{Spec}(A_{\mathfrak{p}})\) Let it be irreducible (resp. that \(A_{p}\) (let it be integral) for all \(p \in U\) ; then \(\text{Spec}(A_{\mathfrak{p}})\) is irreducible (resp. \(A_{p}\) is integral) for all \(p \in \text{Spec}(A)\) .

Let \((\mathrm{Y}_{i})_{i\in\mathrm{I}}\) be the (finite) family of irreducible components of \(\operatorname{Spec}(A)\). Let \(\mathfrak{p}\in U\); since \(\operatorname{Spec}(A_{\mathfrak{p}})\) is irreducible, \(\mathfrak{p}\) contains only one minimal prime ideal of A and therefore belongs to only one of the \(Y_i\) (II, § 4, No. 3, Cor. 2 of Prop. 14). The subsets \(Y_i\cap U\) are closed subsets of U, disjoint, and nonempty since U is dense in \(\operatorname{Spec}(A)\) and irreducible by II, § 4, No. 1, Prop. 7; these are therefore the connected components of U. Their closures \(Y_i\) are the connected components of \(\operatorname{Spec}(A)\) (Cor. 1). This proves that the connected components of \(\operatorname{Spec}(A)\) are irreducible, hence that \(\operatorname{Spec}(A_{\mathfrak{p}})\) is irreducible for every \(\mathfrak{p}\) (No. 8).

Suppose that \(A_{\mathfrak{q}}\) is integral for all \(\mathfrak{q} \in U\). Let \(\mathfrak{p} \in \operatorname{Spec}(A)\). Since \(\operatorname{Spec}(A_{\mathfrak{p}})\) is irreducible, the nilradical of \(A_{\mathfrak{p}}\) is the unique minimal prime ideal of \(A_{\mathfrak{p}}\); it therefore belongs to \(\operatorname{Ass}(A_{\mathfrak{p}})\) (IV, § 1, No. 3, Cor. 1 of Prop. 7), and consequently is equal to \(\mathfrak{q}A_{\mathfrak{p}}\), where \(\mathfrak{q}\) is a prime ideal associated with A (IV, § 1, No. 2, Cor. of Prop. 5). We have \(\mathfrak{q} \in U\) (Remark 1) and \(\mathfrak{q}A_{\mathfrak{q}} \in \operatorname{Ass}(A_{\mathfrak{q}})\) (loc. cit.); since \(A_{\mathfrak{q}}\) is integral, \(\mathfrak{q}\) is zero, therefore \(A_{\mathfrak{p}}\) is an integral domain.

COROLLARY 4.--- Let \(A\) be a Noetherian ring whose spectrum is connected. Suppose there exists an integer \(d \geqslant 1\) such that \(\operatorname{prof}(A_{\mathfrak{p}}) \geqslant 2\) for every prime ideal \(\mathfrak{p}\) of A of height \textgreater{} d.

\begin{enumerate}
\def\labelenumi{\alph{enumi})}
\item
  For every closed subset Z of \(\operatorname{Spec}(\mathbf{A})\) of codimension \(>d\) the space \(\operatorname{Spec}(\mathbf{A}) - \mathbf{Z}\) is connected.
\item
  Let Y and Y' be irreducible components of Spec(A). There exists a sequence \(X_{1}, X_{2}, \ldots, X_{n}\) of irreducible components of Spec(A) such that one has \(X_{1} = Y\) , \(X_{n} = Y'\) and, for \(i = 1, \ldots, n - 1\)
\end{enumerate}

\[
\operatorname{codim} \left(\mathrm{X} _ {i} \cap \mathrm{X} _ {i + 1}, \operatorname{Spec} (\mathrm{A})\right) \leqslant d.
\]

Let \(Z \subset \operatorname{Spec}(A)\) a closed subset of codimension \(>d\) . For every \(\mathfrak{p} \in Z\) , one has \(\dim(A_{\mathfrak{p}}) > d\) , hence \(\operatorname{prof}(A_{\mathfrak{p}}) \geqslant 2\) , which implies that \(\operatorname{prof}_{Z}(A)\) is \(\geqslant 2\) (no. 5, prop. 8) and therefore that \(\operatorname{Spec}(A) - Z\) is connected (Thm. 3).

Let Z denote the union of the sets \(X^{\prime} \cap X^{\prime\prime}\) where \((X^{\prime}, X^{\prime\prime})\) traverses the (finite) set of pairs of irreducible components of Spec(A) such that \(\text{codim}(X^{\prime} \cap X^{\prime\prime}, \text{Spec}(A)) > d\) . By a), the set \(\text{Spec}(A) - Z\) is connected. All irreducible components of \(\text{Spec}(A)\) meet \(\text{Spec}(A) - Z\) ; their traces on \(\text{Spec}(A) - Z\) are the irreducible components of \(\text{Spec}(A) - Z\) (II, § 4, no. 1, prop. 7). Since \(\text{Spec}(A) - Z\) is connected, there exists a sequence \(X_{1}, \ldots, X_{n}\) of irreducible components of \(\text{Spec}(A)\) such that one has \(X_{1} - Z = Y - Z\) , \(X_{n} - Z = Y' - Z\) and \((X_{i} - Z) \cap (X_{i+1} - Z) \neq \emptyset\) for \(1 \leqslant i \leqslant n - 1\) ; in other words, we have \(X_{1} = Y\) , \(X_{n} = Y'\) and \(\text{codim}(X_{i} \cap X_{i+1}) \leqslant d\) .

Remark 2. When A is a Macaulay ring, one may apply the corollary with \(d=1\) ( \(\S\) 2, no. 5).

Example (Complete intersection formed by four coordinate planes of a 4-dimensional affine space).--- Let \(k\) be a field. Let S denote the polynomial ring \(k[\mathrm{T}_{1},\mathrm{T}_{2},\mathrm{T}_{3},\mathrm{T}_{4}]\). Recall (VIII, § 2, No. 4, Th. 3) that every maximal chain of prime ideals of S has length 4. Let \(\mathfrak{m}\) be the ideal of S generated by the \(\mathrm{T}_{i}\), \(\mathfrak{a}\) the ideal generated by \(\mathrm{T}_{1}\mathrm{T}_{2}\) and \(\mathrm{T}_{3}\mathrm{T}_{4}\), and, for \(1 \leqslant i < j \leqslant 4\), let \(\mathfrak{p}_{ij}\) be the ideal generated by \(\mathrm{T}_{i}\) and \(\mathrm{T}_{j}\). The ideals \(\mathfrak{p}_{ij}\) are prime of height 2, their sum is the maximal ideal \(\mathfrak{m}\) and one has \(\mathfrak{a} = \mathfrak{p}_{13} \cap \mathfrak{p}_{14} \cap \mathfrak{p}_{23} \cap \mathfrak{p}_{24}\).

\begin{enumerate}
\def\labelenumi{\alph{enumi})}
\item
  The ring A = S/a is reduced; the irreducible components of Spec(A) are the sets \(X_{ij} = V(\mathfrak{p}_{ij}/\mathfrak{a})\) for i = 1, 2, j = 3, 4, which are of dimension 2 and all contain the point m/a. In particular, Spec(A) is connected and of dimension 2. The intersection of two distinct components \(X_{ij}\) and \(X_{kl}\) is reduced to \(\{m/a\}\) if \(\{i,j\} \cap \{k,l\} = \varnothing\) , of dimension 1 otherwise. It follows that the conclusion of cor. 4 is satisfied with d = 1 (we shall see later that A is a Macaulay ring, so that the hypothesis of corollary 4 is itself satisfied for d = 1).
\item
  Let b be the ideal of S generated by \(T_{1}T_{2}\) , \(T_{1}T_{3}\) , \(T_{2}T_{4}\) , \(T_{3}T_{4}\) , and B the ring S/b. We have \(b = p_{14}p_{23} = p_{14} \cap p_{23}\) . The ring B is reduced. Its spectrum is identified with the closed subset \(X_{14} \cup X_{23}\) of Spec(A); it has two irreducible components (of dimension 2) whose intersection is reduced to a point. The depth of B along this point is strictly positive because B is reduced, and less than 1 by th. 3, hence equal to 1. Thus B is not a Macaulay ring.
\end{enumerate}

\subsection{Depth and normality}

Let A be a Noetherian ring and p a prime ideal of A. We have \(\text{prof}(A_{\mathfrak{p}}) \leqslant \text{ht}(\mathfrak{p})\) (n° 4, cor. 2 of th. 2). If moreover A is reduced, the local ring \(A_{p}\) is reduced (II, § 2, n° 6, prop. 17), which implies:

\begin{enumerate}
\def\labelenumi{\alph{enumi})}
\item
  if ht(p) = 0, A \(_{p}\) is a field;
\item
  if \(\operatorname{ht}(\mathfrak{p}) \geqslant 1\) , one has \(\operatorname{prof}(A_{\mathfrak{p}}) \geqslant 1\) ( \(n^{\circ}\) 1, remark 3).
\end{enumerate}

Conversely:

PROPOSITION 15.--- Let A be a Noetherian ring satisfying the following two conditions:

\begin{enumerate}
\def\labelenumi{(\roman{enumi})}
\item
  for every minimal prime ideal \(\mathfrak{p}\) of A, the ring \(A_{\mathfrak{p}}\) is reduced;
\item
  for every prime ideal \(\mathfrak{p}\) of A of height \(\geqslant 1\) , one has \(\operatorname{prof}(A_{\mathfrak{p}}) \geqslant 1\) . Then A is reduced.
\end{enumerate}

Let \(\mathfrak{n}\) be the nilradical of A. For every minimal prime ideal \(\mathfrak{p}\) of A, we have by (i) \(\mathfrak{n}_{\mathfrak{p}} = 0\) , that is \(\mathfrak{p} \notin \operatorname{Supp}(\mathfrak{n})\) , and a fortiori \(\mathfrak{p} \notin \operatorname{Ass}_{\Lambda}(\mathfrak{n})\) . For every \(\mathfrak{p} \in \operatorname{Spec}(\mathbf{A})\) of height \(\geqslant 1\) , we have by (ii) \(\mathfrak{pA}_{\mathfrak{p}} \notin \operatorname{Ass}_{\mathrm{A}_{\mathfrak{p}}}(\mathrm{A}_{\mathfrak{p}})\) , and a fortiori \(\mathfrak{pA}_{\mathfrak{p}} \notin \operatorname{Ass}_{\mathrm{A}_{\mathfrak{p}}}(\mathfrak{n}_{\mathfrak{p}})\) , whence \(\mathfrak{p} \notin \operatorname{Ass}_{\mathrm{A}}(\mathfrak{n})\) (IV, § 1, no. 2, prop. 5). Thus the set \(\operatorname{Ass}_{\mathrm{A}}(\mathfrak{n})\) is empty, which implies that \(\mathfrak{n}\) is zero.

PROPOSITION 16.--- Let A be a Noetherian integrally closed ring, J an ideal of A of height ≥ 2, and M a reflexive finitely generated A-module. Then prof \(_{A}\) (J; M) ≥ 2.

Let us choose a finite-dimensional vector space V over the field of fractions of A and a lattice N in V isomorphic to M (VII, § 4, No. 2, Remark 1). The associated prime ideals of V/N, being of height 1 (loc. cit., Th. 2), do not contain J; by Remark 2 of No. 1, this implies \(\operatorname{prof}_{\mathrm{A}}(\mathrm{J};\mathrm{V}/\mathrm{N}) \geqslant 1\). On the other hand, the A-module V is divisible and torsion-free, hence injective (A, X, p. 17, Cor. 2 of Prop. 10), which implies \(\operatorname{prof}_{\mathrm{A}}(\mathrm{J};\mathrm{V}) = +\infty\). The inequality \(\operatorname{prof}_{\mathrm{A}}(\mathrm{J};\mathrm{N}) \geqslant 2\) then follows from Prop. 1 of No. 1.

COROLLARY.--- A Noetherian local ring that is integrally closed and of dimension ≥ 2 has depth ≥ 2.

Let A be a normal ring (no. 8) and let p be a prime ideal of A. Then:

\begin{enumerate}
\def\labelenumi{\alph{enumi})}
\item
  if ht(p) = 0, A \(_{p}\) is a field;
\item
  if \(\operatorname{ht}(\mathfrak{p}) = 1\) , \(A_{p}\) is a discrete valuation ring (VII, § 1, no. 7, prop. 11);
\item
  If \(\operatorname{ht}(\mathfrak{p}) \geqslant 2\), one has \(\operatorname{prof}(A_{\mathfrak{p}}) \geqslant 2\) (Corollary of Prop. 16).
\end{enumerate}

Conversely:

THEOREM 4 (Serre).--- Let A be a Noetherian ring satisfying the following conditions:

\begin{enumerate}
\def\labelenumi{(\roman{enumi})}
\item
  for every minimal prime ideal p of A, the ring A \(_{p}\) is reduced ;
\item
  for every prime ideal p of A of height 1, the ring A \(_{p}\) is integrally closed ;
\item
  for every prime ideal \(\mathfrak{p}\) of A of height \(\geqslant 2\) , one has \(\operatorname{prof}(A_{\mathfrak{p}}) \geqslant 2\) .
\end{enumerate}

Then A is normal.

The goal is to prove that the ring \(A_{\mathfrak{p}}\) is integrally closed for every \(\mathfrak{p} \in \operatorname{Spec}(A)\), which we shall do by induction on \(\operatorname{ht}(\mathfrak{p})\). For \(\operatorname{ht}(\mathfrak{p}) \leqslant 1\), this follows from hypotheses (i) and (ii). Suppose therefore that \(\operatorname{ht}(\mathfrak{p}) \geqslant 2\) and that \(A_{\mathfrak{q}}\) is integrally closed for every prime ideal \(\mathfrak{q}\) of A with \(\operatorname{ht}(\mathfrak{q}) < \operatorname{ht}(\mathfrak{p})\). By hypothesis (iii), we have \(\operatorname{prof}(A_{\mathfrak{p}}) \geqslant 2\). By the induction hypothesis and Corollary 3 of Theorem 3 of No. 9, applied to the ring \(A_{\mathfrak{p}}\) and to the closed subset \(\{\mathfrak{p}A_{\mathfrak{p}}\}\) of \(\operatorname{Spec}(A_{\mathfrak{p}})\), the ring \(A_{\mathfrak{p}}\) is an integral domain. Let K be its field of fractions, and let B be a subring of K containing \(A_{\mathfrak{p}}\) and finite over \(A_{\mathfrak{p}}\). It remains to prove that B is equal to \(A_{\mathfrak{p}}\). Let i denote the canonical injection from \(A_{\mathfrak{p}}\) into B. Since B is contained in K, it is an \(A_{\mathfrak{p}}\)-module without torsion, hence of depth \(\geqslant 1\). Moreover, for every prime ideal \(\mathfrak{q}\) of \(A_{\mathfrak{p}}\) distinct from \(\mathfrak{p}A_{\mathfrak{p}}\), the homomorphism \(i_{\mathfrak{q}}: (A_{\mathfrak{p}})_{\mathfrak{q}} \to B_{\mathfrak{q}}\) is bijective since \(A_{\mathfrak{q}}\) is integrally closed. By Lemma 4 of No. 9, applied to the closed subset \(F = \{\mathfrak{p}A_{\mathfrak{p}}\}\) of \(\operatorname{Spec}(A_{\mathfrak{p}})\), the homomorphism i is bijective, which completes the proof of the theorem.

Remark 1.--- A convenient form of Th. 4 is the following: let A be a Noetherian ring such that for every prime ideal p of A the ring A \(_{p}\) is integrally closed or of depth ≥ 2; then A is normal. Indeed, let us verify the hypotheses of Th. 4. Let p ∈ Spec(A). If ht(p) ≤ 1, then we have prof(A \(_{p}\) ) ≤ 1, hence A \(_{p}\) is integrally closed. If ht(p) ≥ 2, the ring A \(_{p}\) is either of depth ≥ 2, or integrally closed, which again implies depth(A \(_{p}\) ) ≥ 2 (no. 1, cor. 2 of prop. 1), whence the desired conclusion. One verifies similarly the following statement: let A be a Noetherian ring such that for every prime ideal p of A, the ring A \(_{p}\) is reduced or of depth ≥ 1; then A is reduced.

COROLLARY 1.--- Let \(A\) be a Noetherian ring, F a closed subset of \(\operatorname{Spec}(A)\), U the complementary open set. We assume that \(\operatorname{prof}_{\mathrm{F}}(\mathrm{A})\) is \(\geqslant 2\) (resp. \(\geqslant 1\) and that, for every \(\mathfrak{p} \in \mathrm{U}\), the ring \(\mathrm{A}_{\mathfrak{p}}\) is integrally closed (resp. reduced). Then A is normal (resp. reduced).

For every \(\mathfrak{p} \in \mathrm{F}\) , one has \(\operatorname{prof}(\mathrm{A}_{\mathfrak{p}}) \geqslant \operatorname{prof}_{\mathrm{F}}(\mathrm{A})\) ( \(n^{\circ}\) 5, prop. 8); it therefore suffices to apply the preceding remark.

COROLLARY 2.-- Let ρ : A → B be a homomorphism of Noetherian rings making B a flat A-module.

\begin{enumerate}
\def\labelenumi{\alph{enumi})}
\item
  If B is normal and faithfully flat over A, then A is normal.
\item
  Suppose that A is normal and that the ring \(\kappa(\mathfrak{p})\otimes_{\mathrm{A}}\mathrm{B}\) is normal when \(\mathfrak{p}\) is a minimal prime ideal of A, and reduced when \(\mathfrak{p}\) is a prime ideal of height 1. Then the ring B is normal.
\item
  Suppose B is normal and faithfully flat over A. Then \(\rho\) is injective (I, § 3, no. 5, prop. 9) and B is reduced, hence A is reduced. Let S be the set of elements of A that are not zero divisors. Since B is flat over A, \(\rho(S)\) is formed of non-zero-divisors in B, so that B is integrally closed in \(\rho(S)^{-1}B\) . Let \(x\) be an element of \(S^{-1}A\) integral over A; then the element \(x \otimes 1_B\) of the ring \(S^{-1}A \otimes_A B\) (which identifies with \(\rho(S)^{-1}B\) ) is integral over B, hence belongs to B. Since B is faithfully flat over A, we have \(x \in A\) (I, § 3, no. 5, prop. 10, (ii)), and A is normal.
\item
  Under the hypotheses of b), it suffices by Remark 1 to prove that for every prime ideal \(\mathfrak{q}\) of B, the local ring \(B_{\mathfrak{q}}\) is normal or has depth \(\geqslant 2\). Let \(\mathfrak{p}\) denote the prime ideal \(\rho^{-1}(\mathfrak{q})\) of A; the local homomorphism \(A_{\mathfrak{p}} \to B_{\mathfrak{q}}\) induced by \(\rho\) makes \(B_{\mathfrak{q}}\) a faithfully flat \(A_{\mathfrak{p}}\)-module (I, § 3, No. 5, Prop. 9). We distinguish four cases:
\end{enumerate}

\begin{enumerate}
\def\labelenumi{\arabic{enumi})}
\item
  If \(\operatorname{ht}(\mathfrak{p}) = 0\), then \(A_{\mathfrak{p}}\) is a field, equal to \(\kappa(\mathfrak{p})\); the ring \(A_{\mathfrak{p}} \otimes_{A} B\) is normal by hypothesis. The same holds for \(B_{\mathfrak{q}}\), which is a ring of fractions thereof.
\item
  If \(\operatorname{ht}(\mathfrak{p}) = 1\) and \(\operatorname{ht}(\mathfrak{q}) \leqslant 1\), then \(A_{\mathfrak{p}}\) is a discrete valuation ring; let \(\pi\) be a uniformizer of \(A_{\mathfrak{p}}\). Since \(B_{\mathfrak{q}}\) is faithfully flat over \(A_{\mathfrak{p}}\), the element \(\pi 1_{B_{\mathfrak{q}}}\) of \(B_{\mathfrak{q}}\) is not a zero divisor, so that the local ring \(B_{\mathfrak{q}}/\pi B_{\mathfrak{q}}\) has dimension 0 (VIII, § 3, No. 1, Cor. 2 of Prop. 1). It is reduced, since it is a ring of fractions of the reduced ring \(\kappa(\mathfrak{p}) \otimes_{A} B\), and consequently it is a field. Therefore \(B_{\mathfrak{q}}\) is a discrete valuation ring (VI, § 1, No. 4, Prop. 2), hence integrally closed.
\item
  \(\operatorname{ht}(\mathfrak{p}) = 1\) and \(\operatorname{ht}(\mathfrak{q}) \geqslant 2\) . Then, by the relation
\end{enumerate}

\[
\dim (B _ {\mathfrak {q}}) = \dim (A _ {\mathfrak {p}}) + \dim (\kappa (\mathfrak {p}) \otimes_ {A} B _ {\mathfrak {q}})
\]

(VIII, § 3, no. 4, cor. 1 of prop. 7), the ring \(\kappa(\mathfrak{p})\otimes_{\mathrm{A}}\mathrm{B}_{\mathfrak{q}}\) has dimension \(\geqslant 1\) . It is reduced by hypothesis, hence of depth \(\geqslant 1\) (no. 1, remark 3). We then have (no. 6, cor. of prop. 11)

\[
\operatorname{prof} \left(\mathrm{B} _ {\mathfrak {q}}\right) = \operatorname{prof} \left(\Lambda_ {\mathfrak {p}}\right) + \operatorname{prof} \left(\kappa (\mathfrak {p}) \otimes_ {\mathrm{A}} \mathrm{B} _ {\mathfrak {q}}\right) \geqslant 1 + 1 = 2.
\]

\begin{enumerate}
\def\labelenumi{\arabic{enumi})}
\setcounter{enumi}{3}
\tightlist
\item
  \(\operatorname{ht}(\mathfrak{p}) \geqslant 2\) . Since \(A_{p}\) is of depth \(\geqslant 2\) (corollary of Prop. 16), the same holds for \(B_{q}\) by loc. cit.
\end{enumerate}

COROLLARY 3.--- Let \(\rho: A \to B\) be a homomorphism of Noetherian rings. Suppose that \(B\) is a flat A-module, that \(A\) is normal, and that \(\kappa(\mathfrak{p}) \otimes_A B\) is normal for all \(\mathfrak{p} \in \operatorname{Spec}(A)\). Then \(B\) is normal.

\section{Macaulay modules and rings}

\subsection{Macaulay modules}

Let A be a Noetherian ring, M a finitely generated A-module, and p a prime ideal of A. If \(p \notin \text{Supp}(M)\) , one has \(M_{p} = 0\) , hence \(\text{prof}_{A_{p}}(M_{p}) = +\infty\) and \(\dim_{A_{p}}(M_{p}) = -\infty\). If \(p \in \text{Supp}(M)\) , one has \(0 \leqslant \text{prof}_{A_{p}}(M_{p}) \leqslant \dim_{A_{p}}(M_{p}) < +\infty\) (§ 1, no. 4, cor. 2 of th. 2).

DEFINITION 1.--- Let A be a Noetherian ring and M a \(\Lambda\) module of finite type. We say that M is Macaulay or is a Macaulay module if, for every maximal ideal \(\mathfrak{m} \in \operatorname{Supp}(M)\) , one has \(\operatorname{prof}_{\mathrm{A}_{\mathfrak{m}}}\left(\mathrm{M}_{\mathfrak{m}}\right) = \dim_{\mathrm{A}_{\mathfrak{m}}}\left(\mathrm{M}_{\mathfrak{m}}\right)\) .

From what precedes, it comes to the same thing to say that one has \(\mathrm{prof}_{\mathrm{A}_{\mathfrak{m}}}\left(\mathrm{M}_{\mathfrak{m}}\right) \geqslant \dim_{\mathrm{A}_{\mathfrak{m}}}\left(\mathrm{M}_{\mathfrak{m}}\right)\) for every maximal ideal \(\mathfrak{m}\) of A. Let A be a Noetherian local ring; for a nonzero finitely generated A-module to be Macaulay, it is necessary and sufficient that its depth be equal to its dimension.

Examples.--- 1) Every A-module of finite length is Macaulay.

\begin{enumerate}
\def\labelenumi{\arabic{enumi})}
\setcounter{enumi}{1}
\tightlist
\item
  Let M' be a direct summand submodule of a finitely generated Macaulay A-module M. Then M' is Macaulay; indeed, for every maximal ideal m of A, the \(A_{m}\) -module \(M'_{m}\) is a direct summand of \(M_{m}\) and we consequently have
\end{enumerate}

\[
\operatorname{prof} _ {A _ {\mathfrak {m}}} \left(M _ {\mathfrak {m}} ^ {\prime}\right) \geqslant \operatorname{prof} _ {A _ {\mathfrak {m}}} \left(M _ {\mathfrak {m}}\right) \geqslant \dim_ {A _ {\mathfrak {m}}} \left(M _ {\mathfrak {m}}\right) \geqslant \dim_ {A _ {\mathfrak {m}}} \left(M _ {\mathfrak {m}} ^ {\prime}\right),
\]

by Remark 4 of § 1, No. 1 and VIII, § 1, No. 4, Prop. 9 c).

\begin{enumerate}
\def\labelenumi{\arabic{enumi})}
\setcounter{enumi}{2}
\tightlist
\item
  Let M be a finitely generated Macaulay module over A and \((x_{1},\ldots,x_{n})\) a M-regular sequence of elements of \(\Lambda\) Then the A-module \(\overline{\mathrm{M}}=\mathrm{M}/(x_{1}\mathrm{M}+\cdots+x_{n}\mathrm{M})\) is Macaulay. Indeed, let m be a maximal ideal of A belonging to the support of \(\overline{M}\) ; one has \(x_{i}\in m\) for every i since \(x_{i}\) vanishes \(\overline{M}\) , and the canonical images of the \(x_{i}\) in \(A_{m}\) form a sequence \(M_{m}\) -regular sequence of elements of \(mA_{m}\) . Consequently (§ 1, no. 4, prop. 7 and VIII, § 3, no. 2, cor. of prop. 3) we have the equalities
\end{enumerate}

\[
\begin{array}{l} \operatorname{prof} _ {\mathrm{A} _ {\mathfrak {m}}} (\overline {{\mathrm{M}}} _ {\mathfrak {m}}) = \operatorname{prof} _ {\mathrm{A} _ {\mathfrak {m}}} (\mathrm{M} _ {\mathfrak {m}}) - n, \\ \dim_ {\mathrm{A} _ {\mathfrak {m}}} (\overline {{\mathrm{M}}} _ {\mathfrak {m}}) = \dim_ {\mathrm{A} _ {\mathfrak {m}}} (\mathrm{M} _ {\mathfrak {m}}) - n, \end{array}
\]

whence our assertion.

\begin{enumerate}
\def\labelenumi{\arabic{enumi})}
\setcounter{enumi}{3}
\tightlist
\item
  Let M be a finitely generated A-module, and a an ideal of A such that \(aM = 0\) . For the A-module M to be Macaulay, it is necessary and sufficient that it be Macaulay as \((A/a)\) -module. Indeed, set B = A/a; let n be a maximal ideal of B and m its inverse image in \(\Lambda\) ; one has \(\dim_{\mathrm{A}_{m}}(\mathrm{M}_{\mathfrak{m}}) = \dim_{\mathrm{B}_{n}}(\mathrm{M}_{\mathfrak{n}})\) and \(\operatorname{prof}_{\mathrm{A}_{m}}(\mathrm{M}_{\mathfrak{m}}) = \operatorname{prof}_{\mathrm{B}_{n}}(\mathrm{M}_{\mathfrak{n}})\) (§ 1, no. 3, cor. of prop. 4).
\end{enumerate}

PROPOSITION 1.--- Let \(\Lambda\) a Noetherian ring, M a finitely generated A-module, \(\mathfrak{p}\) and \(\mathfrak{q}\) of prime ideals of Supp(M) such that \(\mathfrak{p} \subset \mathfrak{q}\) . Suppose \(\dim_{\mathrm{A}_{\mathfrak{q}}}(\mathrm{M}_{\mathfrak{q}}) = \operatorname{prof}_{\mathrm{A}_{\mathfrak{q}}}(\mathrm{M}_{\mathfrak{q}})\) . One then has \(\dim_{\mathrm{A}_{\mathfrak{p}}}(\mathrm{M}_{\mathfrak{p}}) = \operatorname{prof}_{\mathrm{A}_{\mathfrak{p}}}(\mathrm{M}_{\mathfrak{p}})\) and

\[
\dim_ {A _ {q}} (M _ {q}) = \dim_ {A _ {p}} (M _ {p}) + \dim (A _ {q} / p A _ {q}).
\]

This follows directly from cor. 1 of prop. 13 in § 1, no. 7.

COROLLARY.--- Let A be a Noetherian ring and M a finitely generated A-module. The following conditions are equivalent:

\begin{enumerate}
\def\labelenumi{(\roman{enumi})}
\item
  the A-module M is Macaulay
\item
  on \(a \operatorname{prof}_{\mathrm{A}_{\mathfrak{p}}}(\mathrm{M}_{\mathfrak{p}}) = \dim_{\mathrm{A}_{\mathfrak{p}}}(\mathrm{M}_{\mathfrak{p}})\) for every \(\mathfrak{p} \in \operatorname{Supp}(\mathrm{M})\) ;
\item
  we have \(\operatorname{prof}_{\mathrm{F}}(\mathrm{M}) = \operatorname{codim}(\operatorname{Supp}(\mathrm{M}) \cap \mathrm{F}, \operatorname{Supp}(\mathrm{M}))\) for every closed subset \(\mathrm{F}\) of \(\operatorname{Spec}(\mathrm{A})\) ;
\item
  on \(a \operatorname{prof}_{\mathrm{A}}(\mathfrak{p}; \mathrm{M}) = \dim_{\mathrm{A}_{\mathfrak{p}}}(\mathrm{M}_{\mathfrak{p}})\) for every \(\mathfrak{p} \in \operatorname{Supp}(\mathrm{M})\) .
\item
  \(\Rightarrow\) (ii): this follows from Proposition 1.
\item
  \(\Rightarrow\) (iii): by Prop. 8 of § 1, No. 5, \(\operatorname{prof}_{\mathrm{F}}(\mathbf{M})\) is the greatest lower bound of the integers \(\operatorname{prof}(\mathbf{M}_{\mathfrak{p}})\) for \(\mathfrak{p}\) ranging over \(\operatorname{Supp}(\mathbf{M}) \cap \mathbf{F}\). If M is Macaulay, then for such a \(\mathfrak{p}\) we have the equalities \(\operatorname{prof}(\mathbf{M}_{\mathfrak{p}}) = \dim(\mathbf{M}_{\mathfrak{p}}) = \operatorname{codim}(\mathbf{V}(\mathfrak{p}), \operatorname{Supp}(\mathbf{M}))\) (VIII, § 1, No. 4, Prop. 9), whence (iii).
\item
  \(\Rightarrow\) (iv): it suffices to take \(F = V(\mathfrak{p})\) .
\item
  \(\Rightarrow\) (i): this follows from the inequalities \(\operatorname{prof}_{\mathbf{A}}(\mathfrak{p};\mathbf{M})\leqslant\operatorname{prof}(\mathbf{M}_{\mathfrak{p}})\leqslant\dim(\mathbf{M}_{\mathfrak{p}})\) , valid for every \(\mathfrak{p}\in\operatorname{Supp}(\mathbf{M})\) ( \(\S 1,\ n^{\circ}5\) , Remark 3 and \(n^{\circ}4\) , cor. 2 of th. 2).
\end{enumerate}

Remark.--- Let S be a multiplicative subset of A and M a finitely generated Macaulay A-module. Then \(S^{-1}M\) is an \(S^{-1}A\)-module Macaulay. Indeed, let \(\mathfrak{q} \in \operatorname{Spec}(S^{-1}A)\); let us denote \(i_{\mathrm{A}}^{\mathrm{S}} : A \to S^{-1}A\) the canonical homomorphism and \(\mathfrak{p} = (i_{\mathrm{A}}^{\mathrm{S}})^{-1}(\mathfrak{q})\). The ring \((S^{-1}A)_{\mathfrak{q}}\) is identified with \(A_{\mathfrak{p}}\) (II, § 2, No. 5, Prop. 11), and the \(A_{\mathfrak{p}}\)-module \((S^{-1}M)_{\mathfrak{q}}\) with the \(A_{\mathfrak{p}}\)-module \(M_{\mathfrak{p}}\) (II, § 2, No. 7, Prop. 20), which is Macaulay by the corollary.

\subsection{Support of a Macaulay module}

PROPOSITION 2. — Let A be a Noetherian ring and M a finitely generated Macaulay A-module.

\begin{enumerate}
\def\labelenumi{\alph{enumi})}
\item
  The A-module M has no embedded associated prime ideals. \(^{1}\)
\item
  Let X be a closed irreducible subset of Supp(M) and Y a closed subset of X. We have
\end{enumerate}

\[
\operatorname{codim} (Y, X) + \operatorname{codim} (X, \operatorname{Supp} (M)) = \operatorname{codim} (Y, \operatorname{Supp} (M)).
\]

\begin{enumerate}
\def\labelenumi{\alph{enumi})}
\setcounter{enumi}{2}
\item
  The topological space Supp(M) is catenary (VIII, § 1, no. 2, def. 4).
\item
  Let \(X_{1}\) and \(X_{2}\) be irreducible components of Supp(M) and Y a closed subset of \(X_{1} \cap X_{2}\) . We have codim(Y, \(X_{1}\) ) = codim(Y, \(X_{2}\) ).
\item
  Let \(\mathfrak{p} \in \operatorname{Ass}(M)\) . One has \(\operatorname{prof}_{\mathrm{A}}(\mathfrak{p}; M) = 0\) ( \(\S\) 1, \(n^{\circ}\) 1, remark 2), hence \(\dim(M_{\mathfrak{p}}) = 0\) ( \(n^{\circ}\) 1, cor. of prop. 1), which implies that \(\mathfrak{p}\) is a minimal element of \(\operatorname{Supp}(M)\) .
\item
  Suppose first that Y is irreducible. Let p and q be the prime ideals of Supp(M) such that \(Y = V(\mathfrak{q})\) and \(X = V(\mathfrak{p})\) . It follows from prop. 1 that we have
\end{enumerate}

\[
\begin{array}{r l} \operatorname{codim} (Y, X) & = \dim (A _ {q} / p A _ {q}) = \dim (M _ {q}) - \dim (M _ {p}) \\ & = \operatorname{codim} (Y, \operatorname{Supp} (M)) - \operatorname{codim} (X, \operatorname{Supp} (M)). \end{array}
\]

The general case follows from VIII, § 1, no. 2, remark 3.

\begin{enumerate}
\def\labelenumi{\alph{enumi})}
\setcounter{enumi}{2}
\tightlist
\item
  Let X, Y, Z be closed irreducible subsets of Supp(M) such that \(Z \subset Y \subset X\) . The codimension of each of these subsets in Supp(M) is finite (VIII, § 1, no. 4, prop. 9 and § 3, no. 1, cor. 1 of prop. 2). We then deduce from b) the equality
\end{enumerate}

\[
\operatorname{codim} (Z, Y) + \operatorname{codim} (Y, X) = \operatorname{codim} (Z, X)
\]

which implies c) (VIII, § 1, no. 2, prop. 4).

\begin{enumerate}
\def\labelenumi{\alph{enumi})}
\setcounter{enumi}{3}
\tightlist
\item
  By b), we have \(\text{codim}(Y, X_1) = \text{codim}(Y, \text{Supp}(M)) = \text{codim}(Y, X_2)\) .
\end{enumerate}

In particular, if there exists a finitely generated Macaulay A-module M with support equal to Spec(A), the ring A is catenary and consequently every ring of fractions or every quotient ring of A is catenary (VIII, § 1, no. 3, remark 2).

Remark 1. — Under the hypotheses of prop. 2, it can happen that two irreducible components \(X_{1}\) and \(X_{2}\) of Supp(M) have an intersection Y reduced to a point and that we have \(\dim X_{1} \neq \dim X_{2}\) and \(\dim X_{2} \neq \text{codim}(Y, X_{2})\) (cf. exercise 4). However, this cannot happen when the ring A is local, as the following corollary shows.

COROLLARY. — Let A be a Noetherian local ring and M a nonzero finitely generated Macaulay A-module.

\begin{enumerate}
\def\labelenumi{\alph{enumi})}
\item
  All maximal chains of closed irreducible subsets of Supp(M) have length equal to dim(M).
\item
  For every closed subset X of Supp(M), we have
\end{enumerate}

\[
\operatorname{codim} (\mathrm{X}, \operatorname{Supp} (\mathrm{M})) = \dim (\operatorname{Supp} (\mathrm{M})) - \dim (\mathrm{X}).
\]

\begin{enumerate}
\def\labelenumi{\alph{enumi})}
\setcounter{enumi}{2}
\item
  All irreducible components of Supp(M) have the same dimension.
\item
  For every ideal J of A, we have
\end{enumerate}

\[
\operatorname{prof} _ {\mathrm{A}} (\mathrm{J}; \mathrm{M}) = \dim (\mathrm{M}) - \dim (\mathrm{M} / \mathrm{JM}).
\]

\begin{enumerate}
\def\labelenumi{\alph{enumi})}
\tightlist
\item
  A maximal chain of closed irreducible subsets of Supp(M) has as its smallest element \(\{\mathfrak{m}_{\mathrm{A}}\}\) and as its largest element an irreducible component X of
\end{enumerate}

Supp(M). Its length is equal to the codimension of \(\{m_{A}\}\) in X (prop. 2, c)); by prop. 2, b) applied to the closed subsets \(\{m_{A}\}\subset X\) , this is equal to \(\text{codim}(\{\mathfrak{m}_{\Lambda}\},\text{Supp}(M))\) , that is, to \(\dim(M)\) .

\begin{enumerate}
\def\labelenumi{\alph{enumi})}
\setcounter{enumi}{1}
\item
  This is a consequence of a) when the subset X is irreducible (VIII, § 1, no. 2, prop. 5); the general case follows from VIII, § 1, no. 1, prop. 1 and § 1, no. 2, remark 4.
\item
  This is a consequence of b).
\item
  We have \(\operatorname{prof}_{\mathrm{A}}(\mathrm{J};\mathrm{M}) = \operatorname{codim}(\operatorname{Supp}(\mathrm{M})\cap \mathrm{V}(\mathrm{J}),\operatorname{Supp}(\mathrm{M}))\) by the cor. of prop. 1 of no. 1. It suffices then to apply b) with \(\mathrm{X} = \operatorname{Supp}(\mathrm{M})\cap \mathrm{V}(\mathrm{J}) = \operatorname{Supp}(\mathrm{M} / \mathrm{JM})\) (II, § 4, no. 4, cor. of prop. 18).
\end{enumerate}

Remark 2.--- Let M be a finitely generated Macaulay A-module, and let p be an element of Supp(M). In view of th. 2 of § 1, no. 4, we have prof\_A(p; M) \textless{} +∞, and there exists an M-regular sequence of length prof\_A(p; M) consisting of elements of p. Let J denote the ideal of A generated by such a sequence; then the A-module M/JM is Macaulay (no. 1, example 3), and p is a minimal element of its support. Indeed, p contains J and hence belongs to the support of M/JM (II, § 4, no. 4, cor. of prop. 18); by corollary 1 of theorem 2 of § 1, no. 4, the ideal p is contained in an element of Ass(M/JM), but every prime ideal associated with a finitely generated Macaulay module is a minimal element of its support (prop. 2).
